\documentclass[12pt]{article}
\usepackage[english]{babel}

% Commads
\newcommand{\gr}{\selectlanguage{greek}}
\newcommand{\en}{\selectlanguage{english}}
\newcommand{\iso}{\scaleobj{1.3}{\hspace{1pt}{\simeq} \hspace{2pt}}}
\newcommand{\double}{\hspace{2pt}}
\newcommand{\norm}[1]{\left\lVert #1 \right\rVert}
% Dual map of f, written f' so it does not clash with the adjoint f^*.
\newcommand{\dual}[1]{#1'}
% Packages
\usepackage{lmodern}
\usepackage[T1]{fontenc}
\usepackage[utf8]{inputenc}
\usepackage{amsmath,amssymb,amsfonts,amsthm}
\usepackage[alphabetic,nobysame]{amsrefs}
\usepackage[a4paper, margin=1in]{geometry}
\usepackage{MnSymbol}
\usepackage{tikz}
\usepackage{tikz-cd}
\usepackage{mathtools}
\usepackage{hyperref}
\usepackage{physics}
\usepackage{booktabs,array}
\renewcommand{\Tr}{\mathrm{Tr}}
\DeclareMathOperator{\Rank}{rank}
\DeclareMathOperator{\Null}{null}
\DeclareMathOperator{\Ker}{Ker}
\DeclareMathOperator{\adj}{adj}
\usepackage{scalerel}
\usepackage{anyfontsize}
\usepackage{mathrsfs}
% Theorem envs (heading, then a line break)
\newtheoremstyle{breakplain}{\topsep}{\topsep}{\itshape}{0pt}{\bfseries}{.}{\newline}{}
\newtheoremstyle{breakdef}{\topsep}{\topsep}{\normalfont}{0pt}{\bfseries}{.}{\newline}{}
\theoremstyle{breakplain}
\newtheorem{theorem}{Theorem}[section]
\newtheorem{lemma}[theorem]{Lemma}
\newtheorem{corollary}[theorem]{Corollary}
\newtheorem{prop}[theorem]{Proposition}
\theoremstyle{breakdef}
\newtheorem{definition}[theorem]{Definition}
\newtheorem*{example}{Example}
\newtheorem{ex}[subsubsection]{Exercise}
\newtheorem*{ex*}{Exercise}
\theoremstyle{remark}
\newtheorem*{remark}{Remark}
\setlength\parindent{0pt}


% ---------- Cover settings (edit these) ----------
\newcommand{\LeafletTitle}{Linear Algebra}
\newcommand{\LeafletSubtitle}{Theorems and Problems for University Competitions}
\newcommand{\LeafletTags}{SEEMOUS\quad{\Large\textperiodcentered}\quad IMC}
\newcommand{\LeafletDate}{\the\year}
\newcommand{\LeafletAuthor}{}

% ---------- Colors (muted, so the title stays the main thing) ----------
% Hilbert picture: accent (navy) -> plum -> pink -> soft red
\definecolor{plum}{HTML}{6B3A7D}
\definecolor{pink}{HTML}{D9547E}
\definecolor{softred}{HTML}{C0392B}
\definecolor{accent}{HTML}{1F3A5F}  % title

% \CellColor{s}{t}: s = 0..18 position along the palette (navy -> plum -> pink -> red), t = strength in percent
\newcommand{\CellColor}[2]{%
  \ifnum#1<6
    \pgfmathtruncatemacro{\mm}{round(100*(1-#1/6))}%
    \colorlet{cell}{accent!\mm!plum!#2!white}%
  \else\ifnum#1<12
    \pgfmathtruncatemacro{\mm}{round(100*(1-(#1-6)/6))}%
    \colorlet{cell}{plum!\mm!pink!#2!white}%
  \else
    \pgfmathtruncatemacro{\mm}{round(100*(1-(#1-12)/6))}%
    \colorlet{cell}{pink!\mm!softred!#2!white}%
  \fi\fi}

% ---------- Cover ----------
\newcommand{\makecover}{%
  \begin{titlepage}
    \thispagestyle{empty}
    \centering
    \vspace*{1.2cm}

    {\fontsize{45}{44}\selectfont\bfseries\color{accent} \LeafletTitle}\\[0.8em]
    {\Large\itshape \LeafletSubtitle}\\[1.8em]

    {\color{accent}\rule{0.50\textwidth}{0.9pt}}\\[1.8em]

    {\small\scshape\LeafletTags}

    % background illustration: the 10x10 Hilbert matrix, shaded by the entries 1/(i+j-1).
    % Strength follows the entry (it fades to white); the hue runs navy -> plum -> pink -> red along the diagonals.
    \begin{tikzpicture}[remember picture, overlay]
      \begin{scope}[shift={([yshift=-3.0cm]current page.center)}]
        \def\c{0.8}% cell size in cm
        \foreach \i in {1,...,10}{
          \foreach \j in {1,...,10}{
            \pgfmathsetmacro{\v}{1/(\i+\j-1)}
            \pgfmathtruncatemacro{\s}{\i+\j-2}% 0 at the corner ... 18 at the far end
            \pgfmathtruncatemacro{\t}{round(min(70,6+82*\v^0.6+24*sin(10*\s)))}% fade, plus a vibrance bump in the middle
            \ifnum\s<6
              \pgfmathtruncatemacro{\m}{round(100*(1-\s/6))}
              \colorlet{cell}{accent!\m!plum!\t!white}
            \else\ifnum\s<12
              \pgfmathtruncatemacro{\m}{round(100*(1-(\s-6)/6))}
              \colorlet{cell}{plum!\m!pink!\t!white}
            \else
              \pgfmathtruncatemacro{\m}{round(100*(1-(\s-12)/6))}
              \colorlet{cell}{pink!\m!softred!\t!white}
            \fi\fi
            \fill[cell] ({(\j-5.5)*\c-0.5*\c+0.03},{-(\i-5.5)*\c-0.5*\c+0.03})
              rectangle ++({\c-0.06},{\c-0.06});
          }
        }
      \end{scope}
    \end{tikzpicture}

    % border: a hairline frame whose colour runs along the palette (navy -> plum -> pink -> red)
    \begin{tikzpicture}[remember picture, overlay]
      \begin{scope}[shift={(current page.south west)}]
        \pgfmathsetmacro{\fw}{\paperwidth/1cm-1.6}   % frame width in cm
        \pgfmathsetmacro{\fh}{\paperheight/1cm-1.6}  % frame height in cm
        \pgfmathsetmacro{\fl}{0.8}                   % inset from the page edge in cm
        \pgfmathsetmacro{\fr}{\fl+\fw}
        \pgfmathsetmacro{\ft}{\fl+\fh}
        \foreach \k in {0,...,96}{
          \pgfmathtruncatemacro{\st}{round(18*0.5*\k/97)}
          \pgfmathtruncatemacro{\sb}{round(18*(0.5*\k/97+0.5))}
          \pgfmathsetmacro{\xa}{\fl+\fw*\k/97}
          \pgfmathsetmacro{\xb}{\fl+\fw*(\k+1.03)/97}
          \CellColor{\st}{92}
          \draw[cell, line width=1.2pt] (\xa,\ft) -- (\xb,\ft);
          \CellColor{\sb}{92}
          \draw[cell, line width=1.2pt] (\xa,\fl) -- (\xb,\fl);
        }
        \foreach \k in {0,...,139}{
          \pgfmathtruncatemacro{\sl}{round(18*0.5*\k/140)}
          \pgfmathtruncatemacro{\sr}{round(18*(0.5*\k/140+0.5))}
          \pgfmathsetmacro{\ya}{\ft-\fh*\k/140}
          \pgfmathsetmacro{\yb}{\ft-\fh*(\k+1.03)/140}
          \CellColor{\sl}{92}
          \draw[cell, line width=1.2pt] (\fl,\ya) -- (\fl,\yb);
          \CellColor{\sr}{92}
          \draw[cell, line width=1.2pt] (\fr,\ya) -- (\fr,\yb);
        }
      \end{scope}
    \end{tikzpicture}

    \vfill

    \leavevmode{\large \LeafletAuthor}\\[0.4em]
    {\small \LeafletDate}
    \vspace*{1.5cm}
  \end{titlepage}
  \setcounter{page}{1}
}





\begin{document}
\makecover
\tableofcontents
\newpage
\section{Vector spaces}
\subsection{Vector spaces}
\begin{definition}
    Let $\mathbb{F}$ be a field. A \emph{vector space} over $\mathbb{F}$ is a set $V$ together with an addition $+:V\times V\to V$ and a scalar multiplication $\cdot:\mathbb{F}\times V\to V$ such that for all $u,v,w \in V$ and $a,b\in \mathbb{F}$:
    \begin{itemize}
        \item $v+(u+w)=(v+u)+w$.
        \item $v+w=w+v$.
        \item There is an element $0\in V$ with $v+0=v$ for all $v$.
        \item For every $v$ there is an element $-v$ with $v+(-v)=0$.
        \item $a\cdot(b\cdot v)=(ab)\cdot v$.
        \item $1\cdot v=v$.
        \item $a\cdot (v+w)=a\cdot v+a\cdot w$.
        \item $(a+b)\cdot v=a\cdot v+b\cdot v$.
    \end{itemize}
    The elements of $V$ are called \emph{vectors}. Formally the vector space is the triple $(V,+,\cdot)$, but we simply write $V$.
\end{definition}
\begin{example}
    $\mathbb{R}^2=\{(a,b)\double|\double a,b\in \mathbb{R}\}$ is an $\mathbb{R}$-vector space with coordinatewise addition and scalar multiplication.
\end{example}
\begin{example}
    The space $\mathrm{C}[a,b]$ of continuous functions $[a,b]\to\mathbb{R}$ is an $\mathbb{R}$-vector space with pointwise addition and scalar multiplication.
\end{example}
\begin{example}
    $\mathbb{C}$ is both an $\mathbb{R}$-vector space and a $\mathbb{C}$-vector space with the usual operations.
\end{example}
\begin{example}
    The space $\mathbb{C}[x]$ of polynomials with complex coefficients is a vector space over $\mathbb{Q}$, over $\mathbb{R}$ and over $\mathbb{C}$.
\end{example}
\begin{example}
    Let $\mathbb{R}_d[x]=\{p(x)\in\mathbb{R}[x] \double | \double \deg p\leq d \}$ be the polynomials of degree at most $d$. This is an $\mathbb{R}$-vector space. Similarly we write $\mathbb{F}_d[x]$ over any field $\mathbb{F}$.
\end{example}
\begin{example}
    $l^1(\mathbb{R})=\{ (a_n)_{n\geq 0} \subset \mathbb{R} \double | \double \sum_{n=0}^\infty |a_n|< \infty \}$ is an $\mathbb{R}$-vector space.
\end{example}
\begin{example}
    $l^2(\mathbb{R})=\{ (a_n)_{n\geq0}\subset\mathbb{R} \double | \double \sum_{n=0}^\infty |a_n|^2 <\infty \}$ is an $\mathbb{R}$-vector space; closure under addition follows from $|a+b|^2\leq 2|a|^2+2|b|^2$ (or from the Cauchy--Schwarz inequality).
\end{example}
\begin{example}
    The space $c=\{ (a_n)_n \subset \mathbb{R}\double | \double \lim_{n\to \infty}a_n \text{ exists in } \mathbb{R} \}$ of convergent sequences is an $\mathbb{R}$-vector space with termwise addition and scalar multiplication.
\end{example}
\begin{remark}
    In short: \emph{a vector space is an abelian group $(V,+)$ on which the field $\mathbb{F}$ acts}.
\end{remark}
\begin{definition}
    A subset $W$ of a vector space $V$ is a \emph{subspace} if it is a vector space with the operations of $V$. Equivalently, $0\in W$ and $v+aw\in W$ for all $v,w\in W$, $a\in\mathbb{F}$. We write $W\leq V$.
\end{definition}
\begin{example}
    $W_0=\{ f\in \mathrm{C}[a,b] \double | \double f(\frac{a+b}{2})=0\}$ is a subspace of $\mathrm{C}[a,b]$, but $W_1=\{ f\in \mathrm{C}[a,b] \double | \double f(\frac{a+b}{2})=1\}$ is not (it does not contain $0$).
\end{example}
\begin{example}
   $\mathbb{R}_{d-1}[x]$ is a subspace of $\mathbb{R}_d[x]$ and of $\mathbb{R}[x]$, but not of $\mathbb{C}[x]$ or $\mathbb{C}_d[x]$ (as complex vector spaces it is not closed under multiplication by $i$).
\end{example}
\begin{example}
    $V_1=\{ \lambda(1,1) \double | \double \lambda \in \mathbb{R}\}$ and $V_2=\{ \lambda(1,0) \double | \double \lambda \in \mathbb{R}\}$ are subspaces of $\mathbb{R}^2$, but the circle $S_1=\{ (a,b) \double | \double a^2+b^2=1\}$ is not.
\end{example}
\begin{prop}
    Let $(V_i)_{i\in I}$ be a family of subspaces of $V$. Then $\bigcap_{i\in I} V_i$ is a subspace of $V$.
\end{prop}
\begin{remark}
    The index set $I$ may be infinite.
\end{remark}
\begin{definition}
    Vectors $v_1,\ldots, v_n$ are \emph{linearly dependent} if there are $a_1, \ldots , a_n \in \mathbb{F}$, not all zero, with $\sum_{i=1}^na_iv_i=0$.
\end{definition}
\begin{definition}
    Vectors $v_1,\ldots,v_n$ are \emph{linearly independent} if they are not linearly dependent, i.e.\ if $\sum_{i=1}^na_iv_i=0$ implies $a_1=\cdots=a_n=0$.
\end{definition}
\begin{remark}
    Both notions extend to infinite sets: a set is linearly independent if every finite subset is, and linearly dependent otherwise.
\end{remark}
\begin{example}
    In $\mathbb{R}^2$ the sets $\mathcal{B}_1=\{ (1,0), (0,1),(1,1)\}$, $\mathcal{B}_2=\{ (1,0), (0,1),(1,2)\}$ and $\mathcal{B}_3=\{ (0,1),(0,2)\}$ are linearly dependent, while $\mathcal{B}_4=\{ (0,1),(1,1)\}$ and $\mathcal{B}_5=\{ (1,0), (0,2)\}$ are linearly independent.
\end{example}
\begin{example}
    The set $\mathcal{B}_1=\{ x^2+x+1, x^2,1,x+2\}$ is linearly dependent, but $\mathcal{B}_2=\{ x^3,x^2,x+1,2\}$ is linearly independent.
\end{example}
\begin{definition}
    For a subset $\mathcal{B}$ of $V$ the \emph{span} of $\mathcal{B}$ is
    \[\langle \mathcal{B}\rangle_V=\bigcap_{\mathcal{B}\subseteq W \leq V} W,\]
    the smallest subspace of $V$ containing $\mathcal{B}$. We usually drop the subscript.
\end{definition}
\begin{definition}
    For a finite set $\mathcal{B}=\{ v_1,\ldots, v_n \}$ we have $\langle \mathcal{B}\rangle=\{ \sum_{i=1}^na_iv_i \double | \double a_i\in\mathbb{F}\}$. An expression $\sum_{i=1}^na_iv_i$ is called a \emph{linear combination} of $v_1, \ldots, v_n$.
\end{definition}
\begin{example}
    The set $\{ 1,i,2\}\subset\mathbb{C}$ is linearly dependent over both $\mathbb{R}$ and $\mathbb{C}$.
\end{example}
\begin{example}
    The set $\{ 1,i\}\subset\mathbb{C}$ is linearly dependent over $\mathbb{C}$ but not over $\mathbb{R}$.
\end{example}
\begin{remark}
    For an arbitrary subset $\mathcal{B}\subseteq V$,
    \[\langle \mathcal{B}\rangle=\Big\{ \sum_{i=1}^n a_iv_i \double \Big| \double n\in\mathbb{N},\ a_i\in \mathbb{F},\ v_i\in \mathcal{B}\Big\},\]
    the set of all \emph{finite} linear combinations of vectors of $\mathcal{B}$.
\end{remark}
\begin{definition}
    A subset $\mathcal{B}\subseteq V$ is \emph{spanning} if $\langle \mathcal{B}\rangle=V$.
\end{definition}
\begin{example}
    $\mathcal{B}_1=\{ (1,0), (0,1),(1,1)\}$ and $\mathcal{B}_2=\{ (1,0), (0,1),(1,2)\}$ span $\mathbb{R}^2$, but $\mathcal{B}_3=\{ (0,1),(0,2)\}$ does not.
\end{example}
\begin{example}
    The set $\{ x^2+x+1, x^2,1,x+2\}$ spans $\mathbb{R}_2[x]$ but not $\mathbb{R}[x]$.
\end{example}
\begin{example}
    The set $\{ 1,i,2\}$ spans $\mathbb{C}$ over both $\mathbb{R}$ and $\mathbb{C}$.
\end{example}
\begin{example}
    Every set $\mathcal{B}$ spans $\langle \mathcal{B}\rangle$.
\end{example}
\begin{definition}
    A subset that is both spanning and linearly independent is called a \emph{basis}.
\end{definition}
\begin{example}
    $\{(1,0),(0,1)\}$ is a basis of $\mathbb{R}^2$.
\end{example}
\begin{example}
    $\{ 1,i\}$ is a basis of $\mathbb{C}$ over $\mathbb{R}$ but not over $\mathbb{C}$.
\end{example}
\begin{example}
    $\{ x^2,x,1\}$ is a basis of $\mathbb{R}_2[x]$.
\end{example}
\begin{remark}
    A basis in this sense (every vector is a \emph{finite} linear combination) is also called a \emph{Hamel basis}. Using the axiom of choice (Zorn's lemma) one shows that every vector space has a basis, which may be infinite.
\end{remark}
\begin{prop}
    Every linearly independent set extends to a basis.
\end{prop}
\begin{definition}
    All bases of $V$ have the same cardinality, called the \emph{dimension} of $V$ and denoted $\dim_\mathbb{F}(V)$. It need not be finite.
\end{definition}
\begin{remark}
    Over rings instead of fields, the cardinality of a basis need not be an invariant. Rings for which it is are said to have the \emph{Invariant Basis Number} property (IBN).
\end{remark}
\begin{example}
    $\mathbb{R}^2$ and $\mathbb{C}$ have dimension $2$ over $\mathbb{R}$, while $\mathbb{C}$ has dimension $1$ over $\mathbb{C}$.
\end{example}
\begin{example}
    $\mathbb{C}_2[x]$ has dimension $3$ over $\mathbb{C}$ and dimension $6$ over $\mathbb{R}$.
\end{example}
\begin{example}
    $\mathrm{C}[a,b]$ is infinite dimensional over $\mathbb{R}$.
\end{example}
\begin{definition}[*]
    For two vector spaces $V,W$ over the same field, the \emph{external direct sum} (or \emph{direct product}) is
    \[V\times W=\{ (v,w) \double | \double v\in V,\ w\in W\}\]
    with componentwise operations. It contains copies $V\times\{0\}$ and $\{0\}\times W$ of $V$ and $W$. Some authors write $(v,w)$ as a formal combination $v+w$.
\end{definition}
\begin{ex}
    Let $V,W$ be finite dimensional. Find the dimension of $V\times W$.
\end{ex}
\begin{definition}[*]
    For subspaces $V,W$ of $U$ we define their \emph{sum}
    \[V+ W=\{ v+w \double | \double v\in V,\ w\in W \}=\langle V \cup W\rangle_U.\]
\end{definition}
\begin{ex}
    Let $V,W$ be subspaces of a finite dimensional space $U$. Prove the \emph{Grassmann formula}
    \[\dim(V+W)=\dim(V)+\dim(W)-\dim(V\cap W).\]
\end{ex}
\begin{definition}
    If $V,W\leq U$ and $V\cap W=\{0\}$, the sum $V+W$ is called \emph{direct} and is denoted $V\oplus W$.
\end{definition}
\begin{remark}
    The external direct sum $V\times W$ is also commonly denoted $V\oplus W$; we will allow this abuse of notation later on.
\end{remark}
\begin{ex}
    Let $U\leq W$ be finite dimensional subspaces of $V$ of the same dimension. Prove that $U=W$.
\end{ex}

\newpage
\subsection{Linear Functions}
\begin{definition}
     Let $V,W$ be vector spaces over $\mathbb{F}$. A function $f:V\to W$ is \emph{linear} if $f(v+au)=f(v)+af(u)$ for all $u,v \in V$ and $a\in \mathbb{F}$. The set of linear functions $V\to W$ is denoted $\mathcal{L}(V,W)$.
\end{definition}
\begin{prop}
    The zero map is linear, so $\mathcal{L}(V,W)$ is never empty.
\end{prop}
\begin{prop}
    If $g\in\mathcal{L}(U,V)$ and $f\in\mathcal{L}(V,W)$ then $f\circ g\in\mathcal{L}(U,W)$. In particular $\mathcal{L}(V,V)$ is closed under composition.
\end{prop}
\begin{ex}
    Prove that $\mathcal{L}(V,W)$ is an $\mathbb{F}$-vector space under pointwise operations.
\end{ex}
\begin{example}
    The maps $f,g,h:\mathbb{R}^2\to \mathbb{R}^2$ given by $f(x,y)=(y,x)$, $g(x,y)=(x,0)$, $h(x,y)=(2x+3y,x-13y)$ are linear, and so is $(g\circ f)(x,y)=(y,0)$. The map $(x,y)\mapsto (x+x^3-y,y^6+2)$ is not.
\end{example}
\begin{example}
    The maps $f,g,h:\mathbb{R}^2\to \mathbb{R}$ given by $f(x,y)=x-2y$, $g(x,y)=x$, $h(x,y)=y$ are linear, but $x\mapsto x^2$ on $\mathbb{R}$ is not.
\end{example}
\begin{example}
    Differentiation $\mathbb{R}[x]\to \mathbb{R}[x]$, $p(x) \mapsto p'(x)$, is linear.
\end{example}
\begin{example}
    The maps $F,G:\mathrm{C}[0,1]\to \mathbb{R}$ given by $F(f)=\int_0^1 f(x)\, \mathrm{d}x$ and $G(f)=f(1/2)$ are linear, but $f\mapsto \int_0^1 f(x)^2\, \mathrm{d}x$ is not.
\end{example}
\begin{example}
    On the space $c$ of convergent sequences, the Ces\`aro map $(a_n)_n \mapsto \big(\frac{a_1+a_2+\cdots+a_n}{n}\big)_n$ is a linear map $c\to c$.
\end{example}
\begin{example}
    For fixed $h \in \mathcal{L}(V,V)$, the maps $\mathcal{L}(V,V)\to \mathcal{L}(V,V)$ given by $g\mapsto h\circ g$ and $g\mapsto g\circ h$ are linear.
\end{example}
\begin{definition}
    A linear map $f:V\to W$ defines two subspaces, the \emph{kernel} and the \emph{image}:
    \[\Ker(f)=\{x\in V\double | \double f(x)=0 \}\leq V,\qquad \mathrm{Im}(f)=\{f(x) \double | \double x\in V \}\leq W.\]
\end{definition}
\begin{ex}[Categorical definition of the kernel]
    Let $f:V\to W$ be linear and $i:\Ker(f)\to V$ the inclusion. Prove that $f\circ i=0$, and that for every space $U$ and linear $g:U\to V$ with $f \circ g=0$ there is a unique linear $h:U\to\Ker(f)$ with $g=i\circ h$. Show that this property determines $\Ker(f)$ up to unique isomorphism.
    \[\begin{tikzcd}
	{\Ker(f)} & V & W \\
	U
	\arrow["i", from=1-1, to=1-2]
	\arrow["f", from=1-2, to=1-3]
	\arrow["{\exists !h}", dashed, from=2-1, to=1-1]
	\arrow["g"', from=2-1, to=1-2]
    \end{tikzcd}\]
\end{ex}
\begin{example}
    For $f:\mathbb{R}[x]\to \mathbb{R}$, $f(p)=p(1)$, we have $\Ker(f)=\{ (x-1)q(x) \double | \double q\in \mathbb{R}[x] \}$ and $\mathrm{Im}(f)=\mathbb{R}$.
\end{example}
\begin{example}
    For $f:\mathbb{R}^2\to \mathbb{R}^2$, $f(x,y)=(x-2y,0)$, we have $\Ker(f)=\langle (2,1)\rangle$ and $\mathrm{Im}(f)=\langle (1,0)\rangle$.
\end{example}
\begin{ex}
    Let $f:\mathbb{R}[x]\to \mathbb{R}^2$, $f(p)=(p(1),p'(1))$. Describe $\Ker(f)$ explicitly.
\end{ex}
\begin{ex}
    Prove that a linear map is injective if and only if $\Ker(f)=\{0\}$.
\end{ex}
\begin{theorem}[Rank--Nullity]
    Let $f:V\to W$ be linear with $V$ finite dimensional. Then $\dim(V)=\dim\Ker(f)+\dim\mathrm{Im}(f)$.
\end{theorem}
\begin{definition}
    $V$ and $W$ are \emph{isomorphic}, written $V\iso W$, if there is a bijective linear map $V\to W$ (an \emph{isomorphism}).
\end{definition}
\begin{prop}
    Being isomorphic is an equivalence relation on $\mathbb{F}$-vector spaces.
\end{prop}
\begin{prop}
    Let $W,U\leq V$. The map $W\times U\to W+U$, $(w,u)\mapsto w+u$, is linear and surjective, and it is an isomorphism if and only if the sum is direct.
\end{prop}
So a direct sum decomposition breaks a space into smaller pieces that are easier to work with.
\begin{example}
    $\mathbb{R}^2=\langle (1,0)\rangle \oplus\langle (0,1)\rangle\iso \mathbb{R} \times \mathbb{R}$.
\end{example}
\begin{example}
    Let $W=\langle \sin x,\cos x\rangle\leq\mathrm{C}[0,1]$. Then $W\iso \mathbb{R}^2\iso \mathbb{R}\times\mathbb{R}$.
\end{example}
\begin{example}
    Every vector space of finite dimension $d$ over $\mathbb{F}$ is isomorphic to $\mathbb{F}^d$.
\end{example}

\newpage
\subsection{Inner Products}
\begin{definition}
    Let $V$ be a vector space over $\mathbb{F}=\mathbb{R}$ or $\mathbb{C}$. An \emph{inner product} is a map $\langle\cdot,\cdot\rangle:V\times V \to \mathbb{F}$ such that for all $u,v,w\in V$, $a\in\mathbb{F}$:
    \begin{itemize}
        \item $\langle u+av,w\rangle=\langle u,w\rangle+a\langle v,w\rangle$ (linear in the first argument),
        \item $\langle v,w\rangle=\overline{\langle w,v\rangle}$,
        \item $\langle v,v\rangle\geq 0$, with equality if and only if $v=0$.
    \end{itemize}
    Over $\mathbb{R}$ the conjugation does nothing. The standard inner product on $\mathbb{F}^n$ is also written $v\cdot w$.
\end{definition}
\begin{example}
    On $\mathbb{C}^n$, $\langle (v_1,\ldots ,v_n),(w_1,\dots, w_n)\rangle=\sum_{i=1}^nv_i\overline{w_i}$ (the standard inner product).
\end{example}
\begin{example}
    On $\mathrm{C}[0,1]$, $\langle f,g\rangle=\int_0^1 f(x)g(x)\, \mathrm{d}x$.
\end{example}
\begin{example}
    On $l^2(\mathbb{R})$, $\langle (a_n),(b_n)\rangle=\sum_{n=0}^\infty a_nb_n$.
\end{example}
\begin{remark}
    Over general fields one studies bilinear forms instead, and the theory is much richer and stranger; see the section on bilinear forms in the Appendix, and Berger \cites{Berger1,Berger2} and \cite{Kaplansky}.
\end{remark}
\begin{prop}
    If $W\leq V$ and $V$ has an inner product, then its restriction to $W$ is an inner product on $W$.
\end{prop}
\begin{definition}
    Vectors $v,w$ are \emph{orthogonal}, written $v\perp w$, if $\langle v,w\rangle=0$.
\end{definition}
\begin{ex}
    Prove that pairwise orthogonal non-zero vectors are linearly independent.
\end{ex}
\begin{ex}
    Prove that if $\langle x,v\rangle=0$ for all $v\in V$, then $x=0$.
\end{ex}
\begin{definition}
    For $W\leq V$ the \emph{orthogonal complement} of $W$ is $W^{\perp}=\{v\in V\double|\double\langle v,w\rangle=0 \text{ for all } w\in W\}$.
\end{definition}
\begin{prop}
    $W^\perp \leq V$, and if $V$ is finite dimensional then $V=W\oplus W^\perp$.
\end{prop}
\begin{ex}
    Prove that if $W \leq U\leq V$ then $U^\perp \leq W^\perp$.
\end{ex}
\begin{ex}
    Prove that if $V$ is finite dimensional then $(W^\perp)^\perp=W$.
\end{ex}
\begin{prop}
    Let $V$ be finite dimensional and $U,W\leq V$. Then $(U+W)^\perp=U^\perp \cap W^\perp$ and $(U\cap W)^\perp=U^\perp+W^\perp$.
\end{prop}
\begin{definition}
    Vectors $v_1,\dots, v_n$ are \emph{orthonormal} if $\langle  v_i,v_j\rangle=\delta_{i,j}$.
\end{definition}
\begin{definition}
    Let $V$ be a finite dimensional inner product space and $f$ an endomorphism of $V$. There is a unique linear map $f^*:V\to V$ with $\langle f(v),w\rangle=\langle v,f^*(w)\rangle$ for all $v,w$, called the \emph{adjoint} of $f$. If $A$ is the matrix of $f$ in an orthonormal basis, the matrix of $f^*$ in the same basis is $A^*$.
\end{definition}
\begin{theorem}[Gram--Schmidt]
    Let $V$ be an inner product space and $v_1,\dots,v_k$ linearly independent. Then there are orthonormal $w_1,\dots,w_k$ with $\langle w_1,\ldots,w_j\rangle=\langle v_1,\ldots,v_j\rangle$ for every $j$; in particular each $w_i$ is a linear combination of $v_1,\ldots,v_i$.
\end{theorem}
\begin{prop}
    Every finite dimensional inner product space has an orthonormal basis.
\end{prop}
\begin{ex}
    Prove that $\Ker(f^*)=\mathrm{Im}(f)^\perp$ and $\mathrm{Im}(f^*)=\Ker(f)^\perp$.
\end{ex}

\newpage
\section{Matrices}
\begin{definition}
    A matrix over a field $\mathbb{F}$ is a rectangular array of elements of $\mathbb{F}$. The $n\times m$ matrices form the space $\mathbb{M}_{n\times m}(\mathbb{F})$, and we write $\mathbb{M}_n(\mathbb{F})=\mathbb{M}_{n\times n}(\mathbb{F})$. Matrices model linear maps between finite dimensional spaces (see the Appendix).
\end{definition}
\begin{definition}
    For $A=(a_{i,j})\in\mathbb{M}_{n\times m}(\mathbb{F})$ and $B=(b_{i,j})\in\mathbb{M}_{m\times k}(\mathbb{F})$ the \emph{product} $C=AB\in\mathbb{M}_{n\times k}(\mathbb{F})$ is
    \[c_{i,j}=\sum_{r=1}^m a_{i,r}b_{r,j}.\]
\end{definition}
There are other products of matrices, but this one is the most natural and useful; it corresponds to composition of linear maps. Matrices whose sizes do not match in this way cannot be multiplied.
We regard vectors as columns, i.e.\ elements of $\mathbb{M}_{n\times 1}(\mathbb{F})$.
It is often useful to think of the rows or the columns of a matrix as vectors; try to describe the product $AB$ as a way of combining such vectors.
\begin{definition}
    A square matrix $D=(d_{i,j})$ is \emph{diagonal} if $d_{i,j}=0$ whenever $i\neq j$.
\end{definition}
\begin{definition}
    For $A=(a_{i,j})\in\mathbb{M}_{n\times m}(\mathbb{C})$ the \emph{conjugate matrix} is $\overline{A}=(\overline{a_{i,j}})\in\mathbb{M}_{n\times m}(\mathbb{C})$.
\end{definition}
\begin{example}
    If $A=\begin{bmatrix}
        1 & 2 \\
        i & 2-3i
    \end{bmatrix}$ then $\overline{A}=\begin{bmatrix}
        1 & 2 \\
        -i & 2+3i
    \end{bmatrix}$, and if $X=\begin{bmatrix}
        6-7i \\ 3+i \\ 2-i \\ 5
    \end{bmatrix}$ then $\overline{X}=\begin{bmatrix}
        6+7i \\ 3-i \\ 2+i \\ 5
    \end{bmatrix}$.
\end{example}
\begin{definition}
    The \emph{transpose} of $A=(a_{i,j})\in\mathbb{M}_{n\times m}(\mathbb{F})$ is $A^T=(a_{j,i})_{i,j}\in\mathbb{M}_{m\times n}(\mathbb{F})$.
\end{definition}
\begin{example}
    With $A,X$ as above, $A^T=\begin{bmatrix}
        1 & i \\
        2 & 2-3i
    \end{bmatrix}$ and $X^T=\begin{bmatrix}
        6-7i & 3+i & 2-i & 5
    \end{bmatrix}$.
\end{example}
\begin{definition}
    The \emph{conjugate transpose} of $A\in\mathbb{M}_{n\times m}(\mathbb{C})$ is $A^*=\overline{A^T}=\overline{A}^T$.
\end{definition}
\begin{prop}
    Let $A\in \mathbb{M}_{n\times m}(\mathbb{F})$ have rows $v_1,\ldots,v_n$ and $B \in  \mathbb{M}_{m\times k}(\mathbb{F})$ have columns $w_1,\ldots,w_k$:
    \[A = \begin{bmatrix}
    \text{---} & v_1 & \text{---} \\
    \text{---} & v_2 & \text{---} \\
    & \vdots & \\
    \text{---} & v_n & \text{---}
    \end{bmatrix}, \qquad
    B = \begin{bmatrix}
    \mid & \mid & & \mid \\
    w_1 & w_2 & \dots & w_k \\
    \mid & \mid & & \mid
    \end{bmatrix}.\]
    Then
    \[AB = \begin{bmatrix}
    v_1 w_1 & v_1 w_2 & \dots & v_1 w_k \\
    v_2 w_1 & v_2 w_2 & \dots & v_2 w_k \\
    \vdots & \vdots & \ddots & \vdots \\
    v_n w_1 & v_n w_2 & \dots & v_n w_k
    \end{bmatrix}.\]
\end{prop}
\begin{remark}
    We identify $1\times 1$ matrices with elements of the field. For example, if $v=\begin{bmatrix}
       1& 2 & 3
   \end{bmatrix} $ and $w=\begin{bmatrix}
       2 \\1 \\-1
   \end{bmatrix}$ then $vw=[2+2-3]=[1]= 1$.
\end{remark}
\begin{definition}
    $A\in \mathbb{M}_{n}(\mathbb{F})$ is \emph{invertible} if there is $B\in \mathbb{M}_{n}(\mathbb{F})$ with $AB=BA=I_n$. Then $B$ is unique and denoted $A^{-1}$. The invertible $n\times n$ matrices form the group $\mathrm{GL}(n,\mathbb{F})$.
\end{definition}
\begin{remark}
    Non-square matrices cannot have two-sided inverses. This reflects the \emph{Invariant Basis Number} property of fields: $\mathbb{F}^n\iso\mathbb{F}^m$ implies $n=m$ (here $n,m$ may even be cardinals). For more, see \href{https://en.wikipedia.org/wiki/Invariant_basis_number}{[IBN]}.
\end{remark}
\begin{prop}
    Let $A,B\in \mathbb{M}_n(\mathbb{F})$ with $AB=I$. Then $BA=I$, so $B=A^{-1}$.
\end{prop}
\begin{remark}
    This too is special to fields; more precisely it holds over \emph{stably finite} rings, a property slightly stronger than IBN, see \href{https://en.wikipedia.org/wiki/Stably_finite_ring}{[Stably Finite]}. See also \emph{directly finite} (Dedekind-finite) rings, \href{https://en.wikipedia.org/wiki/Dedekind-finite_ring}{[Directly Finite]}.
\end{remark}
\begin{prop}
    For $A\in\mathbb{M}_n(\mathbb{F})$ the following are equivalent:
    \begin{itemize}
        \item $A$ is invertible.
        \item There is no non-zero vector $v$ with $Av=0$.
    \end{itemize}
\end{prop}

Below is a list of special square matrices. We restrict to $A\in\mathbb{M}_n(\mathbb{C})$.
\begin{itemize}
    \item Symmetric: $A=A^T$.
    \item Anti-symmetric (skew-symmetric): $A=-A^T$.
    \item Orthogonal: $A^{-1}=A^T$.
    \item Hermitian: $A=A^*$.
    \item Anti-Hermitian (skew-Hermitian): $A=-A^*$.
    \item Unitary: $A^{-1}=A^*$.
    \item Normal: $AA^*=A^*A$.
    \item Idempotent: $A^2=A$.
    \item Nilpotent: $A^r=O$ for some $r\in \mathbb{N}$.
    \item Involutory: $A^2=I$.
\end{itemize}

\begin{definition}
    $A,B\in\mathbb{M}_n(\mathbb{F})$ are \emph{similar}, written $A\sim B$, if $A=U^{-1}BU$ for some invertible $U$, called a \emph{transition matrix}. If $A,B\in\mathbb{M}_n(\mathbb{C})$ and $U$ can be chosen unitary, $A$ and $B$ are \emph{unitarily similar}.
\end{definition}
\begin{remark}
    Similar matrices represent the same linear map of $V$ with respect to different bases.
\end{remark}
\begin{definition}
    $A,B\in\mathbb{M}_{n\times m}(\mathbb{F})$ are \emph{equivalent} if $A=UBV$ for some invertible $U\in\mathbb{M}_n(\mathbb{F})$, $V\in\mathbb{M}_m(\mathbb{F})$.
\end{definition}
\begin{remark}
    Equivalence makes sense for non-square matrices and is weaker than similarity: it represents the same linear map $V\to W$ with independent changes of basis in $V$ and $W$.
\end{remark}
\begin{definition}
    $A$ is \emph{diagonalizable} if it is similar to a diagonal matrix, and \emph{unitarily diagonalizable} if the transition matrix can be chosen unitary.
\end{definition}
\begin{definition}
    The \emph{commutator} of $A,B\in \mathbb{M}_n(\mathbb{F})$ is $[A,B]=AB-BA$.
\end{definition}
The commutator measures how far $A$ and $B$ are from commuting.
\begin{prop}
    For $A,B,C\in \mathbb{M}_n(\mathbb{F})$ and $x,y \in \mathbb{F}$:
    \begin{itemize}
        \item $[A,A]=O$,
        \item $[A,O]=[A,I]=O$,
        \item $[A,B]=-[B,A]$,
        \item $[A+C,B]=[A,B]+[C,B]$,
        \item $[x A,y B]=xy[A,B]$,
        \item $[A,[B,C]]+[B,[C,A]]+[C,[A,B]]=O$ \textup{(Jacobi identity)}.
    \end{itemize}
\end{prop}
\begin{remark}
    A bilinear, antisymmetric operation satisfying the Jacobi identity is called a \emph{Lie bracket}; these are central in Lie algebras, differential geometry and physics.
\end{remark}

\newpage
\subsection{Exercises}
\begin{ex}
    Show that similarity is an equivalence relation on $\mathbb{M}_n(\mathbb{F})$, which justifies the notation $ A \sim B $.
\end{ex}
\begin{ex}
    Prove or disprove: if $A,B \in \mathrm{GL}(n,\mathbb{C})$ then $AB\in\mathrm{GL}(n,\mathbb{C})$. What about $BA$?
\end{ex}
\begin{ex}
    Prove or disprove: if $A,B \in \mathrm{GL}(n,\mathbb{C})$ then $ A+B\in\mathrm{GL}(n,\mathbb{C}) $.
\end{ex}
\begin{ex}
    Let $A$ be a square matrix with $ A^2=A $. Is $A$ invertible? What if $ A^2=O $?
\end{ex}
\begin{ex}
    Let $A$ be a square matrix with $A^{5} - A^{4} + 4A^{2}+ \pi I=O$. Is $A$ invertible? If so, express $A^{-1}$ as a polynomial in $A$.
\end{ex}
\begin{ex}
    Let $A\sim B$ and suppose $ A^7-3A^4+A^2-2A=O $. Find a polynomial $P$ with $ P(B)=O $.
\end{ex}
\begin{ex}
    Let $A\sim B$ with $A$ invertible. Is $ B $ invertible?
\end{ex}
\begin{ex}
    Let $A,B\in \mathbb{M}_n(\mathbb{C})$ with $AB=BA$ and $B$ invertible. Prove that $AB^{-1}=B^{-1}A$.
\end{ex}
\begin{ex}
    Let $A,B\in \mathbb{M}_n(\mathbb{C})$ with $A\,P(B)=B$ for a polynomial $P$ with $P(0)\neq0$. Prove that $AB=BA$.
\end{ex}
\begin{ex}[\cite{Z3}*{Ex 2.4.9}]
    Assuming that all inverses below exist, show that
    \[(A-B)^{-1}=A^{-1}+A^{-1}(B^{-1}-A^{-1})^{-1}A^{-1},\]
    and derive
    \begin{itemize}
        \item $(I-BA)^{-1}=I+B(I-AB)^{-1}A$,
        \item $(A+I)^{-1}=I-(A^{-1}+I)^{-1}$,
        \item $\det\big((I-A)^{-1}+(I-A^{-1})^{-1}\big)=1$.
    \end{itemize}
\end{ex}
\begin{ex}[\cite{GA17}*{Ex 280}]
    Let $A$ be a square matrix and $ k $ a positive integer with $ kA^{k+1}=(k+1)A^k $. Prove that $ A-I $ is invertible and express its inverse in terms of $A$.
\end{ex}
\begin{ex}[\cite{GA17}*{Ex 281}]
    Let $ A\in\mathbb{M}_n(\mathbb{F}) $ be invertible and $ B=XY $, where $X\in\mathbb{M}_{n\times1}(\mathbb{F})$ and $Y\in\mathbb{M}_{1\times n}(\mathbb{F})$. Prove that $A+B$ is invertible if and only if $a=YA^{-1}X\neq -1$, and that then
    \[(A+B)^{-1}=A^{-1}-\frac{1}{a+1}A^{-1}BA^{-1}.\]
\end{ex}
\begin{ex}
    Let $A,B$ be square matrices with $A+B=AB$. Prove that $AB=BA$.
\end{ex}
\begin{ex}[\cite{GA17}*{Ex 283}]
    Let $A,B$ be square matrices with $ AB-B^2A^2=I $ and $ A^3+B^3=O $. Prove that $ BA-A^2B^2=I $.
\end{ex}
\begin{ex}
    Let $A,B\in \mathbb{M}_n(\mathbb{R})$ be symmetric with non-negative eigenvalues. Prove that $A^3+B^3=(A+B)^3$ if and only if $AB=O$.
\end{ex}

\newpage
\section{Trace and Determinant}
\subsection{Definitions}
\begin{definition}
    The \emph{trace} $\Tr(A)$ of a square matrix $A$ over $\mathbb{F}$ is the sum of its diagonal entries.
\end{definition}
\begin{example}
    If $A=\begin{bmatrix}
        2 & 3\\
        5 & 7
    \end{bmatrix}$ then $\Tr(A)=2+7=9$, and $\Tr(I_3)=1+1+1=3$.
\end{example}
\begin{prop}
    For $A,B \in \mathbb{M}_{n}(\mathbb{F})$ and $\lambda\in\mathbb{F}$:
    \begin{itemize}
        \item $\Tr(A)\in \mathbb{F}$,
        \item $\Tr(A+\lambda B)=\Tr(A)+\lambda \Tr(B)$,
        \item $\Tr(AB)=\Tr(BA)$,
        \item $\Tr(A^T)=\Tr(A)$,
        \item if $A\in\mathbb{M}_{n}(\mathbb{C})$ then $\Tr(\overline{A})=\overline{\Tr(A)}$.
    \end{itemize}
\end{prop}
\begin{definition}
    The \emph{determinant} $\det:\mathbb{M}_n(\mathbb{F})\to\mathbb{F}$ is the unique function that is multilinear and alternating in the columns and satisfies $\det(I_n)=1$.
\end{definition}
Unpacking the definition, with $v,w,v_i$ columns and $\lambda\in\mathbb{F}$:
\begin{itemize}
    \item $\det\begin{bmatrix}
        \vert& \vert& \ldots & \vert\\
        v+w& v_2 &\ldots & v_n\\
        \vert& \vert& \ldots &\vert
    \end{bmatrix} = \det\begin{bmatrix}
        \vert& \vert& \ldots &\vert\\
        v& v_2 &\ldots & v_n\\
        \vert& \vert& \ldots &\vert
    \end{bmatrix} +\det\begin{bmatrix}
        \vert& \vert& \ldots &\vert\\
        w& v_2 &\ldots & v_n\\
        \vert& \vert& \ldots &\vert
    \end{bmatrix}$,
    \item $\det\begin{bmatrix}
        \vert &\vert &\ldots &\vert\\
        v_1& v_2 &\ldots & v_n\\
        \vert &\vert &\ldots &\vert
    \end{bmatrix} = -\det\begin{bmatrix}
        \vert&\vert & \ldots &\vert\\
        v_2& v_1 &\ldots & v_n\\
        \vert&\vert &\ldots &\vert
    \end{bmatrix}$,
    \item $\det\begin{bmatrix}
        \vert& \vert& \ldots & \vert\\
        \lambda v_1& v_2 & \ldots & v_n\\
        \vert& \vert& \ldots & \vert
    \end{bmatrix} = \lambda\det\begin{bmatrix}
        \vert& \vert& \ldots & \vert\\
        v_1& v_2 & \ldots & v_n\\
        \vert& \vert& \ldots & \vert
    \end{bmatrix}$,
\end{itemize}
and the same holds for any other column.
The following properties are the most useful.
\begin{prop}
    For $A,B \in \mathbb{M}_{n}(\mathbb{F})$ and $\lambda\in\mathbb{F}$:
    \begin{itemize}
        \item $\det(A)\in \mathbb{F}$,
        \item $\det(\lambda A)=\lambda^n \det(A)$,
        \item $\det(AB)=\det(A)\det(B)$,
        \item $A \in \mathrm{GL}(n,\mathbb{F}) \iff \det(A)\neq 0$, and then $\det(A^{-1})={1}/{\det(A)}$,
        \item $\det(A^{T})=\det(A)$,
        \item if $A\in \mathbb{M}_{n}(\mathbb{C})$ then $\det(\overline{A})=\overline{\det(A)}$,
        \item $\det(A)=\sum_{\sigma \in S_n}\epsilon(\sigma)\prod_{i=1}^n a_{i,\sigma(i)}$ \textup{(Leibniz formula)}.
    \end{itemize}
\end{prop}
\begin{remark}
    If $A$ has rational (or real, or integer) entries then so do $\Tr(A)$ and $\det(A)$. This gives strong restrictions on such matrices, which we will exploit later.
\end{remark}
\begin{remark}
    Trace and determinant are invariant under similarity: $\Tr(U^{-1}AU)=\Tr(A)$ and $\det(U^{-1}AU)=\det(A)$. Hence they do not depend on the basis, and we can speak of the trace and determinant of an endomorphism.
\end{remark}
\begin{definition}
    For $A\in \mathbb{M}_n(\mathbb{F})$ let $A_{i,j}\in\mathbb{M}_{n-1}(\mathbb{F})$ be the matrix obtained by deleting the $i$-th row and $j$-th column of $A$. The \emph{adjugate} (classical adjoint) of $A$ is
    \[\adj(A)=\big((-1)^{i+j}\det(A_{j,i})\big)_{i,j},\]
    the transpose of the matrix of cofactors.
\end{definition}
\begin{prop}
    For $A,B\in \mathbb{M}_n(\mathbb{F})$ with $n\geq 2$:
    \begin{itemize}
        \item $A\adj(A)=\adj(A)A=\det(A)I$, so $A^{-1} = \frac{1}{\det(A)}\adj(A)$ when $\det(A)\neq 0$,
        \item $\det(\adj(A))=\det(A)^{n-1}$,
        \item $\adj(AB)=\adj(B)\adj(A)$.
    \end{itemize}
\end{prop}
\begin{remark}
    If $A$ has integer entries, so does $\adj(A)$.
\end{remark}
\subsection{Special Determinants}
\begin{prop}[Sarrus]
    For $A=\begin{bmatrix}
        a_1 & a_2 & a_3 \\
        b_1 & b_2 & b_3 \\
        c_1 & c_2 & c_3
    \end{bmatrix}$,
    \[\det(A)=a_1b_2c_3 + a_2b_3c_1+a_3b_1c_2-a_1b_3c_2 - a_2b_1c_3-a_3b_2c_1.\]
\end{prop}
\begin{prop}[Vandermonde determinant]\label{VD}
    \[\det\begin{bmatrix}
        1 & 1 &\cdots & 1 \\
        x_1 & x_2 & \cdots & x_n\\
        x_1^2 & x_2^2 & \cdots & x_n^2 \\
        \vdots & \vdots & \ddots & \vdots \\
        x_1^{n-1} & x_2^{n-1}& \cdots & x_n^{n-1}
    \end{bmatrix}=\prod_{1\leq i < j \leq n} (x_j-x_i).\]
\end{prop}
\begin{prop}[Cauchy determinant]
    Let $A=\big(\frac{1}{x_i-y_j}\big)_{i,j}$ with $x_i\neq y_j$ for all $i,j$. Then
    \[\det(A)=\frac{\prod_{1\leq j<i\leq n}(x_i-x_j)(y_j-y_i)}{\prod_{i=1}^n\prod_{j=1}^n(x_i-y_j)}.\]
\end{prop}
The special case $x_i-y_j=i+j-1$ (e.g.\ $x_i=i$, $y_j=1-j$) is the \emph{Hilbert matrix} $\big(\frac1{i+j-1}\big)_{i,j}$, closely related to Hilbert's inequality for sequences.
\begin{prop}[Circulant determinant]
    A matrix of the form
    \[C=\begin{bmatrix}
        c_0 & c_{n-1} & \cdots & c_2 & c_1\\
        c_1 & c_0 & \cdots & c_3 & c_2\\
        c_2 & c_1 &\cdots & c_4 & c_3\\
        \vdots & \vdots & & \ddots & \vdots \\
        c_{n-1} & c_{n-2} & \cdots & c_1 & c_0
    \end{bmatrix}\]
    is called a \emph{circulant matrix}, and
    \[\det(C)=\prod_{j=0}^{n-1}\big(c_0+c_1 \omega^j+c_2 \omega^{2j}+ \cdots +c_{n-1}\omega^{(n-1)j}\big),\]
    where $\omega=e^{2\pi i/n}$ is a primitive $n$-th root of unity.
\end{prop}
\begin{remark}
    Circulant matrices are diagonalized by the Discrete Fourier Transform (DFT).
\end{remark}
\begin{prop}
    Let $J_n\in\mathbb{M}_n(\mathbb{C})$ be the all-ones matrix and $a,b \in \mathbb{C}$. Then $A=aI_n+bJ_n$, i.e.
    \[A=\begin{bmatrix}
        a+b & b & \cdots & b \\
        b & a+b & \cdots & b \\
        \vdots & & \ddots & \vdots \\
        b & b & \cdots & a+b
    \end{bmatrix},\]
    has $\det(A)=a^{n-1}(a+nb)$.
\end{prop}
\begin{prop}[Smith determinant]
    Let $A=(\gcd(i,j))_{i,j=1}^n$. Then $\det(A)=\prod_{k=1}^n \varphi(k)$, where $\varphi$ is Euler's totient function.
\end{prop}

\newpage
\subsection{Exercises}
\begin{ex}
    Let $A,B$ be square matrices. Prove that $ \Tr(AB)=\Tr(BA) $; hence $[A,B]$ always has trace $0$.
\end{ex}
\begin{ex}[*]
    Let $\mathscr{C}=\langle [A,B] \double | \double A,B \in \mathbb{M}_n(\mathbb{C})\rangle$ be the span of all commutators. Prove that:
    \begin{itemize}
        \item $\mathscr{C}$ is a subspace of $\mathbb{M}_n(\mathbb{C})$; compute its dimension.
        \item $\mathscr{C}=\Tr^{-1}(0)$. (So it is a manifold!)
        \item \textup{(Shoda's theorem)*} Every matrix in $\mathscr{C}$ is a single commutator $[A,B]$.
    \end{itemize}
\end{ex}
\begin{ex}
    Prove that similar matrices have the same trace.
\end{ex}
\begin{ex}
    Let $A,B$ be square matrices. Prove or disprove: $ \Tr(AB)=\Tr(A)\Tr(B) $.
\end{ex}
\begin{ex}
    Prove that there are no $A,B\in \mathbb{M}_n(\mathbb{C})$ with $AB-BA=I$.
\end{ex}
\begin{ex}
    Let $A,B$ be square complex matrices. Prove that $|\Tr(A^*B)|^2 \leq \Tr(A^*A)\Tr(B^*B)$.
\end{ex}
\begin{ex}
    Let $A,B$ be square matrices. Prove or disprove: $ \det(A+B)=\det(A) + \det(B) $.
\end{ex}
\begin{ex}
    Let $ A,B $ be invertible matrices of odd size. Prove that $ AB+BA \neq O $.
\end{ex}
\begin{ex}
    Let $ A,B $ be commuting real matrices. Prove that $ \det(A^2+B^2)\geq 0 $.
\end{ex}
\begin{ex}
    Let $ A,B $ be real matrices of odd size with $ A^2 + B^2 = O $. Prove that $ [A,B]=AB-BA $ is not invertible.
\end{ex}
\begin{ex}
    Prove Proposition \ref{VD}.
\end{ex}
\begin{ex}
    Compute $\det\big(\min(i,j)\big)_{i,j}$ and $\det\big(|i-j|\big)_{i,j}$.
\end{ex}
\begin{ex}
    Let $ A,B,C,D $ be real matrices with $ AC-BD=I $ and $ AD+BC=O $. Prove that
    \begin{itemize}
        \item $ CA-DB=I $ and $ DA+CB=O $,
        \item $ \det(AC)\geq 0 $ and $ (-1)^n\det(BD)\geq 0 $.
    \end{itemize}
\end{ex}
\begin{ex}[\href{https://cdn.artofproblemsolving.com/attachments/6/b/7602596bee912ae3a5085a745ca110424c39d2.pdf}{VRChat Femboy Math Dec 2024}]
    Let $ B\in\mathbb{M}_{2024}(\mathbb{R}) $ be given by $B=(1-\delta_{i,j})_{i,j}$. Compute $\det(B^{-1})$.
\end{ex}
\begin{ex}[\cite{GA17}*{Ex 255}]
    Let $f:\mathbb{M}_{n}(\mathbb{R}) \to \mathbb{R}$ be linear. Prove that $f(X)=\Tr(XC)$ for some $C\in\mathbb{M}_n(\mathbb{R})$, and that if $f(XY)=f(YX)$ for all $X,Y$ then $C=\lambda I$.
\end{ex}
\begin{ex}[*]
    For $\sigma \in S_n$ let $\nu(\sigma)$ be the number of fixed points of $\sigma$. Compute
    \[\sum_{\sigma \in S_n}\frac{\epsilon(\sigma)}{\nu(\sigma)+1}.\]
\end{ex}
\begin{ex}
    Let $A\in \mathbb{M}_n(\mathbb{Z})$ be invertible. Prove that $A^{-1}$ has integer entries if and only if $\det(A)=\pm1$.
\end{ex}
\begin{ex}
    Let $A,P\in \mathbb{M}_n(\mathbb{Z})$ with $\det(P)\neq 0$ and $\det(A)=1$. Prove that some positive power of $B=P^{-1} A P$ has integer entries.
\end{ex}
\begin{ex}
    Let $p$ be a prime and $X_0,\ldots,X_{p-1}$ independent random variables with $X_i\sim \mathrm{Ber}(1/2)$, i.e.\ each is $0$ or $1$ with probability $1/2$. What is the probability that the circulant matrix
    \[\begin{bmatrix}
    X_0 & X_1 & \cdots & X_{p-1}\\
    X_{p-1} & X_0 & \cdots & X_{p-2}\\
    \vdots & \vdots & \ddots & \vdots \\
    X_1 & X_2 & \cdots & X_0
    \end{bmatrix}\]
    is invertible?
\end{ex}

\newpage
\section{Rank inequalities and subspace inclusions}
This section is devoted to rank inequalities, which come from inclusions between subspaces naturally attached to a map, mainly kernels and images.
\begin{definition}
    The \emph{rank} of a matrix $A$ is $\Rank(A)=\dim\mathrm{Im}(A)$, and its \emph{nullity} is $\Null(A)=\dim\Ker(A)$.
\end{definition}
\begin{example}
    $A=\begin{bmatrix}
        1 & 0 & 0\\0 & 1 & 0\\ 0& 0 & 0
    \end{bmatrix}$ has $\Rank(A)=2$ and $\Null(A)=1$, and so does $B=\begin{bmatrix}
        4 & 3 & 7\\0 & 2 & 10\\ 0& 0 & 0
    \end{bmatrix}$.
\end{example}
\begin{prop}
    For $A\in \mathbb{M}_{n\times m}(\mathbb{F})$, $\Rank(A)+\Null(A)=m$.
\end{prop}
\begin{prop}
    For $A\in \mathbb{M}_{n\times m}(\mathbb{F})$, $\Rank(A)\leq \min(n,m)$.
\end{prop}
\begin{prop}
    The rank of $A$ equals the largest size of a square submatrix of $A$ with non-zero determinant. In particular, if $\Rank(A)=k$ then every $(k+1)\times(k+1)$ submatrix has determinant $0$.
\end{prop}
\begin{prop}
    For $A,B\in\mathbb{M}_{n}(\mathbb{F})$, $\Rank(AB)\leq \min(\Rank(A),\Rank(B))$. If one of them is invertible, equality holds; the converse is false.
\end{prop}
\begin{remark}
    Rank is invariant under similarity, so we can speak of the rank and nullity of an endomorphism.
\end{remark}
\begin{prop}
    For $A,B \in\mathbb{M}_{n\times k}(\mathbb{F})$, $\Ker(A)\cap \Ker(B)\subseteq \Ker(A+B)$.
\end{prop}
\begin{prop}
    For $A,B\in\mathbb{M}_{n}(\mathbb{F})$, $\Rank(A+B)\leq \Rank(A)+\Rank(B)$.
\end{prop}
\begin{theorem}[Frobenius' inequality]
    For matrices $A,B,C$ such that $ABC$ is defined,
    \[\Rank(ABC)+\Rank(B)\geq \Rank(AB)+\Rank(BC).\]
\end{theorem}
\begin{corollary}[Sylvester's inequality]
    For $A\in\mathbb{M}_{n\times k}(\mathbb{F})$ and $B\in\mathbb{M}_{k\times m}(\mathbb{F})$,
    \[\Rank(AB)+k\geq \Rank(A)+\Rank(B).\]
\end{corollary}
\begin{corollary}[Sylvester's inequality, kernel form]
    For $A,B\in\mathbb{M}_{n}(\mathbb{F})$, $\Null(AB)\leq \Null(A)+\Null(B)$.
\end{corollary}
Sylvester's inequality is the case $B=I_k$ of Frobenius' inequality (with $A,B,C$ renamed), and in contests it is the version used most often. Both generalize to more matrices, but such generalizations rarely appear in contests; the interested reader can find more in \cite{T}.
\begin{prop}
    If $A\in \mathbb{M}_n(\mathbb{C})$ satisfies $A^2=A$, then $\Tr(A)=\Rank(A)$.
\end{prop}
\begin{prop}
    For $A,B\in \mathbb{M}_n(\mathbb{F})$, $\Rank(AB)=\Rank(B)-\dim(\Ker(A)\cap \mathrm{Im}(B))$.
\end{prop}
\begin{prop}
    For $A,B\in\mathbb{M}_{n}(\mathbb{F})$, $\Rank(AB-I)=\Rank(BA-I)$.
\end{prop}
\begin{proof}
    Equivalently, $\Null(AB-I)=\Null(BA-I)$. If $x\in \Ker(AB-I)$ then $ABx=x$, so $BA(Bx)=Bx$, i.e.\ $Bx\in \Ker(BA-I)$. Moreover $B$ is injective on $\Ker(AB-I)$: if $Bx=0$ then $x=ABx=0$. Hence $B$ maps $\Ker(AB-I)$ injectively into $\Ker(BA-I)$ and $\Null(AB-I)\leq \Null(BA-I)$. Exchanging the roles of $A$ and $B$ gives the reverse inequality.
\end{proof}
\begin{prop}
    For every $A\in\mathbb{M}_{n\times m}(\mathbb{C})$,
    \[\Rank(A)=\Rank(A^*)=\Rank(\overline{A})=\Rank(AA^*)=\Rank(A^*A).\]
\end{prop}

\begin{example}
    Let $A,B\in \mathbb{M}_n(\mathbb{C})$ with $\Rank(A)+\Rank(B)=n$, such that $A+B$ is invertible and commutes with $A$ (or with $B$). Prove that for every $k\in \mathbb{N}$:
    \begin{itemize}
        \item $A^k+B^k$ is invertible,
        \item $\Tr((A+B)^k)=\Tr(A^k)+\Tr(B^k)$.
    \end{itemize}
    \begin{proof}
        From $\Rank(A)+\Rank(B)=n$ we get $\Null(A)+\Null(B)=n$. Moreover $\Ker(A)\cap \Ker(B)\subseteq\Ker(A+B)=\{ 0\}$ because $A+B$ is invertible, so by the Grassmann formula
        \[V=\Ker(A) \oplus \Ker(B).\]
        If $A+B$ commutes with $A$ then $A^2+BA=A^2+AB$, so $AB=BA$ (similarly if it commutes with $B$). Now we show $AB=O$. Write any $v$ as $v=v_A+v_B$ with $v_A\in\Ker(A)$, $v_B\in\Ker(B)$. Then
        \[ABv=ABv_A+ABv_B=BAv_A+A(Bv_B)=0.\]
        So $AB=BA=O$, hence $(A+B)^k=A^k+B^k$ for every $k$, and both claims follow: $A^k+B^k=(A+B)^k$ is invertible, and taking traces gives the second.
    \end{proof}
\end{example}
\begin{theorem}
    Let $A\in \mathbb{M}_n(\mathbb{C})$ and $P$ a polynomial with $P(A)=O$. If $P=RQ$ with $\gcd(R,Q)=1$, then $V=\Ker(Q(A)) \oplus \Ker(R(A))$.
\end{theorem}

\newpage
\subsection{On rank one matrices}
Matrices of rank one appear often in competitions, and they are mostly handled with the following facts.
\begin{prop}
    Every rank one matrix can be written as $vw^T$ for some non-zero column vectors $v,w$, and then $\Tr(vw^T)=w^Tv$.
\end{prop}
\begin{prop}
    Every rank one matrix $A$ satisfies $A^2=\Tr(A) A$. Moreover:
    \begin{itemize}
        \item if $\Tr(A)=0$, then $A$ is nilpotent and similar to $E_{1,2}$, the matrix with a single $1$ in position $(1,2)$ and zeros elsewhere (in a suitable basis it sends $e_2\mapsto e_1$ and every other $e_i\mapsto 0$);
        \item if $\Tr(A)\neq0$, then $A$ is diagonalizable with the single non-zero eigenvalue $\Tr(A)$.
    \end{itemize}
\end{prop}
\begin{remark}
    The matrix $vw^T$ is how the tensor product $v\otimes w$ looks in coordinates.
\end{remark}
\subsection{Exercises}
\begin{ex}
    Let $A\in \mathbb{M}_{n}(\mathbb{C})$ with $A^2=O$. Prove that $\Rank(A)\leq \frac{n}{2}$.
\end{ex}
\begin{ex}
    Let $A\in \mathbb{M}_{n}(\mathbb{C})$ with $A^2-3A+2I=O$. Prove that $\Rank(A-I)+\Rank(A-2I)=n$.
\end{ex}
\begin{ex}[Gazeta Matematic\u{a} 465]
    Let $A\in \mathbb{M}_{n\times m}(\mathbb{C})$ and $B\in \mathbb{M}_{m\times n}(\mathbb{C})$.
    Prove that $\Rank(I_n-AB)-\Rank(I_m-BA)=n-m$.
\end{ex}
\begin{ex}[Gazeta Matematic\u{a} 537]
    Let $A,B\in \mathbb{M}_n(\mathbb{C})$ with $A^2=A$ and $B^2=B$. Prove that
    \[\mathrm{Im}(AB-BA)=\mathrm{Im}(A+B-I)\cap \mathrm{Im}(A-B).\]
\end{ex}
\begin{ex}
    Let $A\in \mathbb{M}_n(\mathbb{C})$. Prove that $\Ker(A-I)\subseteq \mathrm{Im}(A)$.
\end{ex}
\begin{ex}
    Let $A\in \mathbb{M}_n(\mathbb{C})$ with $A^2=A$. Prove that $\Rank(A-I)=n-\Tr(A)$.
\end{ex}
\begin{ex}
    Let $A,B\in \mathbb{M}_n(\mathbb{C})$. Show that $\Ker(A)\cap \Ker(B)\subseteq \Ker(A+B)$.
\end{ex}
\begin{ex}
    Let $A,B\in \mathbb{M}_n(\mathbb{C})$. Show that $\Ker(A^*A)=\Ker(A)$, $\Ker(AA^*)=\Ker(A^*)$ and $\Ker(A^*A+B^*B)=\Ker(A)\cap \Ker(B)$.
    Does the analogous statement hold for $\Ker(AA^*-BB^*)$?
\end{ex}
\begin{remark}
    Here you may use that $x^*C^*Cx=\norm{Cx}^2$, so $x^*C^*Cx=0$ implies $Cx=0$, for any matrix $C$ and vector $x$.
\end{remark}
\begin{ex}
    Let $A,B\in \mathbb{M}_n(\mathbb{C})$ with $A^*B=O$. Prove that $\Rank(A^*A+B^*B)\leq \Rank(AA^*+BB^*)$.
\end{ex}
\begin{ex}[SEEMOUS 2020 P3]
    Let $n$ be a positive integer, $ k\in \mathbb{C}$ and $A\in \mathbb{M}_n(\mathbb{C})$ with $ \Tr(A) \neq 0 $ and
    \[\Rank(A)+\Rank(\Tr(A)\cdot I-kA)=n.\]
    Find $\Rank(A)$.
\end{ex}
\begin{ex}
    Let $A,B\in \mathbb{M}_n(\mathbb{C})$ with $\Rank(A)+\Rank(B)=\Rank(A+B)$ and $AB+BA=O$. Show that $AB=O$.
\end{ex}
\begin{ex}
    Let $A,B\in \mathbb{M}_n(\mathbb{C})$ be symmetric with $(AB+BA-A^2-I)^2=AB^2-B^2A$. Find $\Rank(A^2+B^2)$.
\end{ex}
\begin{ex}
    Let $A,B\in \mathbb{M}_n(\mathbb{C})$ have the same rank and let $k\geq 1$ be an integer. Prove that $A^{k+1}B^k=A$ if and only if $B^{k+1}A^k=B$.
\end{ex}
\begin{ex}
    Let $A,B\in \mathbb{M}_n(\mathbb{C})$. Prove that the following are equivalent:
    \begin{itemize}
        \item $AB=A$ and $\Rank(I-A)+\Rank(B)=n$;
        \item $BA=B$ and $\Rank(I-B)+\Rank(A)=n$.
    \end{itemize}
\end{ex}
\begin{ex}
    Let $A_1,\ldots, A_k\in \mathbb{M}_n(\mathbb{C})$ with $A_i^2=A_i$ for all $i$. Prove that
    \[\Rank\Big(I-\prod_{i=1}^{k} A_i\Big)\leq \sum_{i=1}^k\Null(A_i).\]
\end{ex}
\subsection*{Rank one}
\begin{ex}
    Let $A,B\in \mathbb{M}_n(\mathbb{C})$ with $\Rank(B)=1$. Prove that $\det(A+xB)=\det(A)+ax$ for some $a\in \mathbb{C}$ and all $x\in\mathbb{C}$.
    If $A,B$ are rational (or real), must $a$ be rational (or real)?
\end{ex}
\begin{ex}
    Let $A,X\in \mathbb{M}_n(\mathbb{R})$ with $\Rank(X)=1$. Prove that
    \[\det(A+X)\det(A-X)\leq\det(A)^2.\]
\end{ex}
\begin{ex}
    Let $A,B\in \mathbb{M}_n(\mathbb{R})$ with $\Rank(AB-BA+I)=1$. Prove that
    \[\Tr(ABAB)-\Tr(A^2B^2)=\frac{n(n-1)}{2}.\]
\end{ex}
\begin{ex}
    Let $A,B\in \mathbb{M}_n(\mathbb{C})$ with $\Rank([A,B])=1$. Prove that $[A,B]^2=O$.
\end{ex}
\begin{ex}
    Let $A\in \mathbb{M}_n(\mathbb{C})$ have rank one. Prove that $\det(A+I)=\Tr(A)+1$.
\end{ex}
\begin{ex}
    Let $A,B\in \mathbb{M}_n(\mathbb{C})$ with $\Rank(AB+BA)=1$. Prove that $[A,B]$ is not invertible.
\end{ex}
\begin{ex}
    Let $A,B\in \mathbb{M}_n(\mathbb{C})$ with $A^2+B^2=2AB$. Prove that $[A,B]$ is not invertible, and that if $\Rank(A-B)=1$ then $[A,B]=O$.
\end{ex}
\begin{ex}
    Let $A,B\in \mathbb{M}_n(\mathbb{R})$ be such that $X=A+iB$ has rank one. Prove that $\Tr(A^2+B^2)\geq 0$, and that if $\Tr(A^2+B^2)=  0$ then $A^2+B^2=O$.
\end{ex}
\begin{ex}
    Let $A,B\in \mathbb{M}_n(\mathbb{C})$ with $A^2=[A,B]$ and $\Rank(A+B)=1$. Prove that $A^3=ABA=O$.
\end{ex}
\begin{ex}
    Let $A,B\in \mathbb{M}_n(\mathbb{C})$ with $A^2=[A,B]$ and $\Rank(A-B)=1$. Prove that $ABA=O$.
\end{ex}
This is essentially the previous exercise, and was posed at SEEMOUS.

\newpage
\section{Block Matrices}
Block matrices are included because many rank arguments are most naturally written in block form. Kernel arguments give perfectly good proofs as well, but they sometimes require quotient spaces, and block matrices give quicker and cleaner answers.
\begin{definition}
    A square matrix of size $n$ is written in \emph{block form} as
    \[M=\begin{bmatrix}
        X&Y\\
        Z&W
    \end{bmatrix}\]
    with $X\in \mathbb{M}_{l\times k}(\mathbb{F})$, $Y\in \mathbb{M}_{l\times (n-k)}(\mathbb{F})$, $Z\in \mathbb{M}_{(n-l)\times k}(\mathbb{F})$ and $W\in \mathbb{M}_{(n-l)\times (n-k)}(\mathbb{F})$.
\end{definition}
\begin{prop}
    Let $M=\begin{bmatrix}
        M_{1,1}&M_{1,2}\\
        M_{2,1}&M_{2,2}
    \end{bmatrix}$ and $N=\begin{bmatrix}
        N_{1,1}&N_{1,2}\\
        N_{2,1}&N_{2,2}
    \end{bmatrix}$ be square matrices with compatible block decompositions. Then $M+N$ and $MN$ have the same block form and
    \[(M+N)_{i,j}=M_{i,j}+N_{i,j},\qquad (MN)_{i,j}=\sum_k M_{i,k}N_{k,j}.\]
    The order of the factors in $M_{i,k}N_{k,j}$ matters, since matrices do not commute in general.
\end{prop}
\begin{example}
    Let $A\in \mathbb{M}_{2 \times 4}(\mathbb{C})$ and $B\in \mathbb{M}_{4 \times 2}(\mathbb{C})$ with
    \[BA=\begin{bmatrix}
        1 & 0 &-1 & 0\\
        0 & 1 & 0 & -1\\
        -1& 0 & 1 & 0 \\
        0 & -1 & 0 & 1
    \end{bmatrix}.\]
    Compute $AB$.
\end{example}
\begin{proof}
    Write $A=\begin{bmatrix}
        A_1 & A_2
    \end{bmatrix}$ and $B=\begin{bmatrix}
        B_1 \\ B_2
    \end{bmatrix}$ with $A_i,B_i\in \mathbb{M}_{2}(\mathbb{C})$. Then
    \[BA=\begin{bmatrix}
        B_1A_1 & B_1A_2\\
        B_2A_1 & B_2A_2
    \end{bmatrix}=\begin{bmatrix}
        I & -I\\
        -I & I
    \end{bmatrix},\]
    so $B_1A_1=I$, $B_1A_2=-I$ and $B_2A_1=-I$. Hence $A_1=B_1^{-1}$, $A_2=-B_1^{-1}=-A_1$ and $B_2=-A_1^{-1}=-B_1$. Therefore $AB=A_1B_1+A_2B_2=I+(-A_1)(-B_1)=2I$.
\end{proof}
\begin{definition}
    For $A\in\mathbb{M}_{n\times m}(\mathbb{F})$ and $B\in\mathbb{M}_{k\times l}(\mathbb{F})$ let
    \[A \oplus B=\begin{bmatrix}
        A & O_{n \times l}\\
        O_{k \times m} & B
    \end{bmatrix}\in\mathbb{M}_{(n+k)\times(m+l)}(\mathbb{F}),\]
    and define $\bigoplus_{i=1}^n M_i$ inductively.
\end{definition}
\begin{prop}
    Let $A=\begin{bmatrix}
        X & Y \\
        O & Z
    \end{bmatrix}$ with $X$ and $Z$ square, possibly of different sizes. Then $\det(A)=\det(X)\det(Z)$, and the same holds for $\begin{bmatrix}
        X & O \\
        Y & Z
    \end{bmatrix}$.
\end{prop}
This is the block version of the fact that the determinant of a triangular matrix is the product of its diagonal entries, and by induction it extends to block triangular matrices with more blocks.
\begin{prop}
    If $A=\begin{bmatrix}
        X & Y \\
        O & Z
    \end{bmatrix}$ then $\Rank(A)\geq \Rank(X)+\Rank(Z)$, with equality if $Y=O$.
\end{prop}
\begin{theorem}[Schur complement]
    Let $M=\begin{bmatrix}
        X & Y \\
        Z & W
    \end{bmatrix}$ with $X$ square and invertible, and let $S=W-ZX^{-1}Y$ be the \emph{Schur complement} of $X$ in $M$. Then
    \[M=\begin{bmatrix}
        I & O\\
        ZX^{-1} & I
    \end{bmatrix}\begin{bmatrix}
        X & O\\
        O & S
    \end{bmatrix}\begin{bmatrix}
        I & X^{-1}Y\\
        O & I
    \end{bmatrix}.\]
    In particular $\det(M)=\det(X)\det(S)$ and $\Rank(M)=\Rank(X)+\Rank(S)$.
\end{theorem}
\begin{prop}
    Let $A \in \mathbb{M}_{n \times m}(\mathbb{F})$ and $B \in \mathbb{M}_{m \times n}(\mathbb{F})$. Then $\det(I_n -AB)=\det(I_m-BA)$.
\end{prop}
\begin{example}
    For every $A\in\mathbb{M}_n(\mathbb{F})$, $\Rank(A^2-A)+n=\Rank(A)+\Rank(A-I)$.
\end{example}
\begin{proof}
    Consider $A\oplus (A-I)$, whose rank is $\Rank(A)+\Rank(A-I)$. Multiplying by invertible block matrices such as $\begin{bmatrix} I & O\\
        I & I
    \end{bmatrix}$ does not change the rank, so we may perform block row and column operations. Be careful: adding a \emph{left} multiple of a block row to another, or a \emph{right} multiple of a block column, is allowed, but the multiplier must be on the correct side. Always write down the invertible matrices you are multiplying by.
    \[\begin{bmatrix}
    A & O\\
    O & A-I
    \end{bmatrix}\to
    \begin{bmatrix}
    A & O\\
    A & A-I
    \end{bmatrix}\to
    \begin{bmatrix}
    A & O\\
    A & I-A
    \end{bmatrix}\to
    \begin{bmatrix}
    A & A\\
    A & I
    \end{bmatrix}\to
    \begin{bmatrix}
    A-A^2 & O\\
    A & I
    \end{bmatrix}\to
    \begin{bmatrix}
    A-A^2 & O\\
    O & I
    \end{bmatrix}.\]
    The last matrix has rank $\Rank(A-A^2)+n$.
\end{proof}
\begin{remark}
    See also \cite{Pop2020}.
\end{remark}
With this in mind, try to solve some of the rank exercises of the previous section using block matrices.

\newpage
\subsection{Exercises}
\begin{ex}[\cite{Z} Ex 2.64]
    Let $A\in\mathbb{M}_{m\times n}(\mathbb{F})$ and $B\in\mathbb{M}_{n\times p}(\mathbb{F})$. Show that
    \[\Rank\begin{bmatrix}
    O & A\\
    B& I_n
    \end{bmatrix}=n+\Rank(AB).\]
\end{ex}
\begin{ex}[\cite{Z} Ex 2.45]
    Prove or disprove:
    \begin{itemize}
        \item $A^2=B^2 \implies A=B$ or $A=-B$;
        \item $\Rank(A)=\Rank(B)\implies \Rank(A^2)=\Rank(B^2)$;
        \item $\Rank(A+kB)\leq \Rank(A)+k\Rank(B)$ for every scalar $k$;
        \item $\Rank(A-B)\leq \Rank(A) - \Rank(B)$;
        \item $\Rank(AB)=0 \implies \Rank(BA)=0$;
        \item $\Rank(AB)=0 \implies \Rank(A)\Rank(B)=0 $.
    \end{itemize}
\end{ex}
\begin{ex}
    Let $A\in \mathbb{M}_n(\mathbb{C})$ with $A^3=A$. Prove that $\Rank(A)+\Rank(A-I)+\Rank(A+I)=2n$.
\end{ex}
\begin{ex}
    Let $A\in\mathbb{M}_{n\times m}(\mathbb{F})$ and $B\in\mathbb{M}_{m\times n}(\mathbb{F})$. Prove that $ \det(I_n-AB)=\det(I_m-BA)$.
    (Hint: compute the determinant of $M=\begin{bmatrix}
        I_n & A \\ B & I_m
    \end{bmatrix}$ in two ways.)
\end{ex}
\begin{ex}
    Let $A,B\in \mathbb{M}_n(\mathbb{C})$. Prove that $\Rank(A-ABA)-\Rank(B-BAB)=\Rank(A)-\Rank(B)$.
\end{ex}
\begin{ex}
    Let $A,B\in \mathbb{M}_n(\mathbb{C})$ with $A^2=A$ and $B^2=B$. Prove that $\Rank(A-BA)-\Rank(B-AB)=\Rank(A)-\Rank(B)$.
    (Hint: $\Rank(XY)=\,?$)
\end{ex}
\begin{ex}
    Let $A,B \in \mathbb{M}_n(\mathbb{F})$. Prove that $\det\begin{bmatrix}
        A & B \\
        B & A
    \end{bmatrix}=\det(A+B)\det(A-B)$.
\end{ex}
\begin{ex}
    Let $A,B \in \mathbb{M}_n(\mathbb{R})$. Prove that $\det\begin{bmatrix}
        A & -B \\
        B & A
    \end{bmatrix}=|\det(A+iB)|^2\geq 0$.
\end{ex}
\begin{ex}
    Let $A,B \in \mathbb{M}_n(\mathbb{F})$. Prove that $\Rank\begin{bmatrix}
        A & I \\
        I & B
    \end{bmatrix}=n+\Rank(I-AB)$.
\end{ex}
\begin{ex}
    Prove that $\Rank\begin{bmatrix}
        A & AB \\
        CA & CAB
    \end{bmatrix}=\Rank(A)$.
\end{ex}

\newpage
\section{Eigenvectors and Eigenvalues}
\begin{definition}
    Let $f$ be an endomorphism of $V$ (with matrix $A$). A subspace $W\leq V$ is \emph{invariant} under $f$ if $f(W)\subseteq W$.
\end{definition}
To simplify computations we try to write $V$ as a direct sum of invariant subspaces $W_i$; then to compute $Aw$ for $w\in W_i$ we only need to understand $A$ on $W_i$, not on all of $V$.
The following proposition makes this precise.
\begin{prop}
    Let $A\in\mathbb{M}_n(\mathbb{F})$ and $V=\bigoplus_i W_i$ with each $W_i$ invariant under $A$. Choosing a basis of each $W_i$ and putting them together gives a basis of $V$ with transition matrix $U$ such that
    \[A =UMU^{-1},\qquad M=\bigoplus_i M_i,\]
    where $M_i$ is the matrix of $A|_{W_i}$.
\end{prop}
So $f$ is the ``sum'' of the maps $f|_{W_i}$, each acting on its own invariant subspace.
\begin{definition}
    An \emph{eigenvector} of an endomorphism $f$ is a non-zero vector $v$ with $f(v)=\lambda v$ for some $ \lambda \in \mathbb{F} $; such $ \lambda $ are the \emph{eigenvalues} of $f$.
\end{definition}
An eigenvector spans a one-dimensional invariant subspace.
\begin{prop}
    Eigenvectors with distinct eigenvalues are linearly independent. In particular, an $n\times n$ matrix with $n$ distinct eigenvalues is diagonalizable.
\end{prop}
A square matrix is diagonalizable exactly when $V$ is a direct sum of one-dimensional invariant subspaces, i.e.\ when $V$ has a basis of eigenvectors. Note also that an $n\times n$ matrix has at most $n$ eigenvalues. This is often used to pin down the eigenvalues: for example, if one shows that $\lambda$ an eigenvalue implies $\lambda^2$ an eigenvalue, then the eigenvalues lie among $0$ and roots of unity.
\begin{example}
    Let $F,G$ be endomorphisms of $\mathbb{R}^n$ and $a\neq 0$ a real number with $[F,G]=F\circ G-G\circ F=aF$. Show that
    \begin{itemize}
        \item $F^k\circ G- G\circ F^k=akF^k$ for all $k\geq 1$;
        \item $F^k=O$ for some $k$.
    \end{itemize}
\end{example}
\begin{proof}
    The first claim is by induction on $k$, the case $k=1$ being the hypothesis. If it holds for $k$, multiply it on the right by $F$:
    \[F^k\circ G \circ F- G\circ F^{k+1}=akF^{k+1}.\]
    Since $G\circ F=F\circ G-aF$, the first term equals $F^k\circ(F\circ G-aF)=F^{k+1}\circ G-aF^{k+1}$, and substituting gives $F^{k+1}\circ G-G\circ F^{k+1}=a(k+1)F^{k+1}$.

    For the second claim consider the linear map $L(X)=[X,G]$ on the finite dimensional space $\mathcal{L}(\mathbb{R}^n,\mathbb{R}^n)$. By the first claim $L(F^k)=akF^k$, so whenever $F^k\neq O$ the number $ak$ is an eigenvalue of $L$. These are pairwise distinct, and $L$ has only finitely many eigenvalues, so $F^k=O$ for some $k$.
\end{proof}
The same argument shows that $[A,B]=A$ forces $A$ to be nilpotent.

\subsection{Characteristic and minimal polynomial}
The eigenvectors for $\lambda$ are the non-zero vectors of $\Ker(A-\lambda I)$, since $Av=\lambda v \iff (A-\lambda I)v=0$; and this kernel is non-trivial exactly when $\det(A-\lambda I)=0$. This leads to the following definition.
\begin{definition}
    For $A\in \mathbb{M}_n(\mathbb{F})$ the \emph{characteristic polynomial} is $\chi_A(x)=\det(A-xI)$, a polynomial of degree $n$ whose roots in $\mathbb{F}$ are exactly the eigenvalues of $A$.
\end{definition}
\begin{definition}
    For $A\in \mathbb{M}_n(\mathbb{C})$ the \emph{spectrum} $\sigma(A)$ is the set of eigenvalues of $A$. Since they are the roots of $\chi_A$, $\sigma(A)$ is non-empty and has at most $n$ elements.
\end{definition}
\begin{theorem}[Spectral mapping]
    For $A\in \mathbb{M}_n(\mathbb{C})$ and $f\in \mathbb{C}[x]$, $\sigma(f(A))=f(\sigma(A))$.
\end{theorem}
\begin{theorem}[Cayley--Hamilton]
    Every $A\in \mathbb{M}_n(\mathbb{F})$ satisfies $\chi_A(A)=O$.
\end{theorem}
Some useful properties of the characteristic polynomial:
\begin{itemize}
    \item if $A$ has entries in $\mathbb{F}$ then $ \chi_A(x)\in \mathbb{F}[x] $;
    \item $ \chi_A(0)=\det(A) $;
    \item if $A\sim B$ then $ \chi_A=\chi_B $;
    \item if $\chi_A$ has no repeated roots then $A$ is diagonalizable.
\end{itemize}
\begin{remark}
    For $ A \in \mathbb{M}_{n}(\mathbb{F})$ consider $ I_A=\{P\in \mathbb{F}[x] \double|\double P(A)=O\} $. It contains $\chi_A$ by Cayley--Hamilton, and it is an ideal of $\mathbb{F}[x]$. Since $\mathbb{F}[x]$ is a principal ideal domain, $I_A=\langle m_A\rangle$ for a unique monic polynomial $m_A$, the non-zero element of $I_A$ of least degree. It is called \emph{the minimal polynomial} of $A$.
\end{remark}
Some properties of the minimal polynomial:
\begin{itemize}
    \item if $P(A)=O$ then $m_A\mid P$; in particular $m_A\mid \chi_A$;
    \item every eigenvalue of $A$ is a root of $m_A$;
    \item $A$ is diagonalizable if and only if $m_A$ splits with no repeated roots; in particular, if $m_A$ has a repeated root then $A$ is not diagonalizable;
    \item $m_A\in \mathbb{F}[x]$.
\end{itemize}
The last property means that if $A$ is real or rational, so is its minimal polynomial; we will see applications in the exercises.

\newpage
\subsection{Jordan Canonical Form}
Not every matrix is diagonalizable, i.e.\ there may not be enough eigenvectors to form a basis. The Jordan form describes exactly what happens instead.
\begin{definition}
    The \emph{Jordan block} of size $m$ and eigenvalue $a$ is $N_m(a)=aI_m + N_m(0)$, where $N_m(0)=(\delta_{i+1,j})_{i,j}$ has ones directly above the diagonal and zeros elsewhere.
\end{definition}
\begin{example}
    \[N_2(3)=\begin{bmatrix}3 & 1 \\
    0 & 3
    \end{bmatrix},\qquad
    N_3(2)=\begin{bmatrix}
        2 &1&0\\
        0 &2 &1 \\
        0 & 0 &2
    \end{bmatrix},\qquad
    N_4(e)=\begin{bmatrix} e & 1 & 0 & 0\\
        0 &e & 1 & 0 \\
        0 & 0 & e & 1  \\
       0 & 0 &0 &e
    \end{bmatrix}.\]
\end{example}
\begin{definition}
    A \emph{generalized eigenvector} of degree $m$ of $A$, for the eigenvalue $\lambda$, is a vector $v$ with $(A-\lambda I)^mv=0$ but $(A-\lambda I)^{m-1}v\neq 0$.
\end{definition}
\begin{theorem}[Jordan canonical form]
    Let $A\in\mathbb{M}_n(\mathbb{C})$ (or any $A$ whose characteristic polynomial splits) with minimal polynomial $m_A(x)=\prod_i(x-a_i)^{n_i}$, the $a_i$ distinct. Then $A$ is similar to a block diagonal matrix
    \[J=\bigoplus_k N_{m_k}(\lambda_k)\]
    of Jordan blocks, unique up to the order of the blocks. Each $\lambda_k$ is one of the $a_i$, and for each $a_i$ the largest block with eigenvalue $a_i$ has size exactly $n_i$.
\end{theorem}
So $V$ is a direct sum of invariant subspaces, one per block, each spanned by a chain of generalized eigenvectors of degrees $1,2,\ldots,m_k$.
This may look abstract, but it gives a concrete way to think about matrices. A typical application: to prove a statement for all matrices, when the statement is preserved by similarity and behaves well with block diagonal matrices, it is enough to prove it for a single Jordan block.

\begin{definition}
    Let $A\in \mathbb{M}_n(\mathbb{C})$ and $\lambda\in\sigma(A)$. The \emph{geometric multiplicity} of $\lambda$ is the number of Jordan blocks with eigenvalue $\lambda$, which equals $\dim\Ker(A-\lambda I)$. The \emph{algebraic multiplicity} of $\lambda$ is the sum of the sizes of these blocks, which equals its multiplicity as a root of $\chi_A$.
\end{definition}
\par\medskip\noindent
\begin{minipage}{0.58\textwidth}
\resizebox{\linewidth}{!}{
$J = \left[\begin{array}{ccc|cc|cc|cc|ccc|c}
3&1& & & & & & & & & & & \\
 &3&1& & & & & & & & & & \\
 & &3& & & & & & & & & & \\ \hline
 & & &3&1& & & & & & & & \\
 & & & &3& & & & & & & & \\ \hline
 & & & & &3&1& & & & & & \\
 & & & & & &3& & & & & & \\ \hline
 & & & & & & &2&1& & & & \\
 & & & & & & & &2& & & & \\ \hline
 & & & & & & & & &2&1& & \\
 & & & & & & & & & &2&1& \\
 & & & & & & & & & & &2& \\ \hline
 & & & & & & & & & & & &4
\end{array}\right]$}
\end{minipage}
\hfill
\begin{minipage}{0.38\textwidth}
Writing $\mathscr{A}$ for the algebraic and $\mathscr{G}$ for the geometric multiplicity, here $\mathscr{A}(3)=7$, $\mathscr{G}(3)=3$, $\mathscr{A}(2)=5$, $\mathscr{G}(2)=2$ and $\mathscr{A}(4)=\mathscr{G}(4)=1$.
\end{minipage}
\par\medskip
The Jordan form also gives a non-algebraic description of the minimal polynomial: it is the $\mathrm{lcm}$ of the characteristic polynomials of the Jordan blocks. In the example above, $m_J(x)=(x-3)^3(x-2)^3(x-4)$.

\newpage
\subsection{Exercises}
\begin{ex}
    Prove that $AB$ and $BA$ have the same characteristic polynomial.
\end{ex}
\begin{ex}
    Find $A,B\in\mathbb{M}_n(\mathbb{C})$ such that $AB$ and $BA$ have:
    \begin{itemize}
        \item Different ranks;
        \item Different Jordan forms but the same minimal polynomial.
    \end{itemize}
\end{ex}
\begin{ex}
   Let $A\in\mathbb{M}_n(\mathbb{C})$. Prove that $\det(A)$ is the product of the eigenvalues (counting multiplicity) of $A$ and $\tr(A)$ is their sum.
\end{ex}
\begin{ex}
    Let $A\in\mathbb{M}_n(\mathbb{C})$. Apply the vieta formulas on the characteristic polynomial of $A$. Prove the identity $B^2=\tr(B)B-I\det(B)$ for $B\in \mathbb{M}_2(\mathbb{C})$.
\end{ex}
\begin{ex}
    Let $p$ be prime. Can there be a non zero matrix $A$ in $\mathbb{M}_{p-1}(\mathbb{Q})$ with $A^p=A$ and $\tr(A)=0$?
\end{ex}
\begin{ex}\label{P}
    Let $p$ be an odd prime. Let $A\in\mathbb{M}_p(\mathbb{Q})$ such that $\det(A^p+I)=0$ and $\det(A+I)\neq 0$. Prove that $\tr(A)$ is an eigenvalue of $A+I$.
\end{ex}
\begin{ex}
    Let $A,B \in \mathbb{M}_n(\mathbb{C})$ have the same characteristic polynomial. Must they be similar? What if they have the same minimal polynomial?
\end{ex}
\begin{ex}
    Is there a real square matrix $M$ of odd size with $M^2+I=O$? What about $M^2+M+I=O$?
\end{ex}
\begin{ex}
    Prove the proposition on rank one matrices given in the section on rank.
\end{ex}
\begin{ex}
    Let $A\in \mathbb{M}_n(\mathbb{C})$ with $A^2=A$. Prove that $\Rank(A)=\Tr(A)$.
\end{ex}
\begin{ex}
    Let $\mathscr{G}$ be a finite multiplicative subgroup of $\mathrm{GL}(n,\mathbb{C})$ and $S=\sum_{M\in \mathscr{G}} M$. Prove that if $\Tr(S)=0$ then $S=O$.
\end{ex}
\begin{ex}
    Let $A,B \in \mathbb{M}_n(\mathbb{C})$ with $AB-BA=A$. Prove that all eigenvalues of $A$ are $0$.
    What do its characteristic and minimal polynomials look like?
    (Hint: Newton's identities.)
\end{ex}
\begin{ex}[SEEMOUS 2023 P1]
    Let $A,B \in \mathbb{M}_n(\mathbb{C})$ with $A=AB-BA+A^2B-2ABA+BA^2+A^2BA-ABA^2$. Prove that $\det(A)=0$.
\end{ex}
\begin{ex}
    Let $A\in \mathbb{M}_n(\mathbb{C})$. Prove that $A^T \sim A$.
\end{ex}
\begin{ex}
    Let $A,B\in \mathbb{M}_n(\mathbb{C})$ and $f:\mathbb{M}_n(\mathbb{C}) \to \mathbb{M}_n(\mathbb{C})$, $f(X)=AX-XB$. Prove that $f$ is an isomorphism if and only if $A$ and $B$ have no common eigenvalue.
\end{ex}
\begin{ex}
    Use the Jordan form to describe the matrices of size $2,3,4,5$ annihilated by $x^3+1$, and those annihilated by $x^4+1$.
\end{ex}
\begin{ex}
    Let $A\in \mathbb{M}_n(\mathbb{C})$ with $A^m=I$, and let $\zeta$ be a primitive $m$-th root of unity.
    Prove that $\Rank(A-I)+\Rank(A-\zeta I)+\cdots+\Rank(A-\zeta^{m-1}I)=n(m-1)$.
\end{ex}
\begin{ex}
    Let $A,B,C,D\in \mathbb{M}_n(\mathbb{C})$ with $\det(AC)\neq0$ and $A^kB=C^kD$ for all $k\geq 1$. Prove that $B=D$.
\end{ex}
\begin{ex}
    Let $A$ be a matrix with non-zero trace. Prove that $A^2=\Tr(A) A$ if and only if $\Rank(A)=1$.
\end{ex}
\begin{ex}
    Let $A\in \mathbb{M}_n(\mathbb{C})$. Prove that
    \[\det(A)=\frac{1}{n!}\det\begin{bmatrix}
    \Tr(A) & 1 & 0 & \cdots & \cdots & 0 \\
    \Tr(A^2) & \Tr(A) & 2 & 0 & \cdots & 0 \\
    \Tr(A^3) & \Tr(A^2) & \Tr(A) & 3 & \cdots & 0 \\
    \vdots & \vdots & \vdots & \vdots & \ddots & \vdots \\
    \Tr(A^{n-1}) & \Tr(A^{n-2}) & \cdots & \cdots & \Tr(A) & n-1 \\
     \Tr(A^{n}) & \Tr(A^{n-1}) & \cdots & \cdots & \Tr(A^2) & \Tr(A)
    \end{bmatrix}.\]
\end{ex}
This is a consequence of Newton's identities, and it shows that if $\Tr(A^i)=0$ for $i=1,\ldots ,n$ then $\det(A)=0$.
\begin{ex}[SEEMOUS TST AUTH 2025]
    Let $A,B \in \mathbb{M}_n(\mathbb{C})$ with $A+I$ invertible and $AB+A=B^2+B$, and let $k$ be the algebraic multiplicity of the eigenvalue $-1$ of $B$. Prove that $\Rank(A-B)=k$.
\end{ex}
\begin{ex}[SEEMOUS TST AUTH 2021]
    Let $A\in \mathbb{M}_n(\mathbb{C})$ have minimal polynomial of degree $m$, and let $B=\big(\Tr(A^{i+j-2})\big)_{i,j=1}^m\in \mathbb{M}_m(\mathbb{C})$ (with $A^0=I$). Prove that $A$ is diagonalizable if and only if $B$ is invertible.
\end{ex}
\begin{ex}
    Let
    \[A=\begin{bmatrix}
        a & -a & -1 & 0 \\
        a & -a & 0 & -1 \\
        1 & 0 & a & -a \\
        0 & 1 & a & -a
    \end{bmatrix}.\]
    Find all rational $a$ for which $A=X^2$ for some matrix $X$ with rational entries.
\end{ex}
\begin{ex}
    Let $A,B\in\mathbb{M}_n(\mathbb{C})$, $n\geq2$, with
    \begin{itemize}
        \item $A^2=B^2$,
        \item $A^3+BAB=2I_n$.
    \end{itemize}
    Prove that:
    \begin{itemize}
        \item $A^3=I_n$;
        \item if $A$ and $B$ have no common eigenvalue, then $A+B=O$.
    \end{itemize}
\end{ex}

\newpage
\section{Diagonal forms}
\begin{theorem}[Schur's theorem]
    For every $A\in \mathbb{M}_n(\mathbb{C})$ there is a unitary $U$ such that $A=U^*TU$ with $T$ upper triangular; the diagonal entries of $T$ are the eigenvalues of $A$ (with algebraic multiplicity).
\end{theorem}
\begin{remark}
    One can equally make $T$ lower triangular. This is helpful when working with products such as $AA^*$, where the diagonal entries are easier to compute for one triangular shape than for the other.
\end{remark}
\begin{theorem}[Simultaneous diagonalization]
    If $A_1,\ldots,A_k$ are pairwise commuting diagonalizable matrices, there is an invertible $V$ such that every $V^{-1} A_i V$ is diagonal.
\end{theorem}
\begin{theorem}[Simultaneous triangularization]
    If $A_1,\ldots,A_k\in\mathbb{M}_n(\mathbb{C})$ pairwise commute, there is a unitary $U$ such that every $U^{*} A_i U$ is upper triangular.
\end{theorem}
\begin{theorem}[Simultaneous unitary diagonalization]
    If $A_1,\ldots,A_k\in\mathbb{M}_n(\mathbb{C})$ are pairwise commuting and each is unitarily diagonalizable (i.e.\ normal), there is a unitary $U$ such that every $U^{*} A_i U$ is diagonal.
\end{theorem}
It is worth taking a minute to play with eigenvectors under the assumption of commutativity: if $AB=BA$, then $B$ maps each eigenspace of $A$ into itself.
\begin{prop}[Laffey's theorem*]
    If $\Rank([A,B])=1$ then $A$ and $B$ are simultaneously triangularizable, though not necessarily by a unitary matrix.
\end{prop}
This does not generalize directly to more than two matrices and is included only for completeness.

\subsection{Exercises}
\begin{ex}
    Let $A\in \mathbb{M}_n(\mathbb{C})$. Prove that there is a unitary $U$ such that $A=U^*LU$ with $L$ lower triangular, whose diagonal entries are the eigenvalues of $A$.
\end{ex}

\newpage
\section{Transpose of A Matrix and Inner Products}
\begin{theorem}[Spectral theorem]
    Let $A\in\mathbb{M}_n(\mathbb{C})$ be normal. Then $A$ is unitarily diagonalizable: there is $U$ with $U^*U=I$ and $U^*AU=D$, where $D$ is diagonal with the eigenvalues of $A$ on the diagonal.
\end{theorem}
\begin{prop}
    For $A\in \mathbb{M}_n(\mathbb{C})$ the matrix $A^*A$ is Hermitian, hence unitarily diagonalizable, and its eigenvalues are real and non-negative: if $A^*Ax=\lambda x$ with $x\neq0$, then $\norm{Ax}^2=x^*A^*Ax=\lambda\norm{x}^2$, so $\lambda = \norm{Ax}^2/\norm{x}^2\geq0$.
\end{prop}

\begin{prop}\label{prop:KA}
    Let $A\in\mathbb{M}_n(\mathbb{C})$ with $\Ker(A^*)=\Ker(A)$. Then there is a unitary matrix $U$ such that
    \[A=U^*\begin{bmatrix}
        P & O\\
        O & O
    \end{bmatrix}U\]
    with $P$ square and invertible.
\end{prop}
\begin{proof}
    Recall $\Ker(A^*)=\mathrm{Im}(A)^\perp$. Hence $\Ker(A)=\mathrm{Im}(A)^\perp$, so
    \[\mathbb{C}^n=\mathrm{Im}(A) \oplus \Ker(A)\]
    is an orthogonal decomposition. Choose an orthonormal basis of $\mathrm{Im}(A)$ and one of $\Ker(A)$; together they form an orthonormal basis of $\mathbb{C}^n$, and the transition matrix to it is unitary. Both subspaces are invariant under $A$, and $A$ is zero on $\Ker(A)$, so in the new basis $A$ has the form $P\oplus O$, where $P$ is the matrix of $A|_{\mathrm{Im}(A)}$. Finally $P$ is invertible: if $w\in\mathrm{Im}(A)$ and $Aw=0$, then $w\in \mathrm{Im}(A)\cap\Ker(A)=\{0\}$.
\end{proof}

\subsection{Exercises}
\begin{ex}
    Let $A\in \mathbb{M}_n(\mathbb{R})$ with $A+A^T=I$. Compute $\det(A)$.
\end{ex}
\begin{ex}
    Let $A\in \mathbb{M}_n(\mathbb{C})$. Using eigenvalues, prove that if $\Tr(A^*A)=0$ then $A=O$.
\end{ex}
\begin{ex}
    Find all $A\in \mathbb{M}_n(\mathbb{C})$ with $A^{2023}=AA^*-A^*A$.
\end{ex}
\begin{ex}
    Let $A\in \mathbb{M}_n(\mathbb{C})$ with $(AA^*)^2=A^*A$. Prove that $A^*A=AA^*$.
\end{ex}
\begin{ex}
    Let $\omega$ be a non-real cube root of unity. Is there a matrix $A\in \mathbb{M}_{2025}(\mathbb{C})$ with $\Tr(A)=0$ and
    \[AA^*+\omega A^*-\omega^2A=2I?\]
\end{ex}
\begin{ex}[SEEMOUS 2025 P3]\label{S25P3}
    Let $A\in \mathbb{M}_n(\mathbb{C})$ with $A^*A^2=AA^*$. Prove that $A^2=A$.
\end{ex}
\begin{ex}
    Let $P\in \mathbb{C}[x]$. Prove that if $A^*=P(A)$ then $A$ is unitarily diagonalizable.
\end{ex}
\begin{ex}
    Let $A\in \mathbb{M}_n(\mathbb{C})$ have all entries in $\{1,-1\}$.
    Prove that if $A$ is normal, then it has an eigenvalue $\lambda$ with $|\lambda|\geq \sqrt{n}$. Does this hold without the normality assumption?
\end{ex}
\begin{ex}[SEEMOUS 2014 P3]
    Let $A\in \mathbb{M}_n(\mathbb{C})$ and $\alpha\in \mathbb{C}$ with $A-A^*=2\alpha I$. Prove that:
    \begin{itemize}
        \item $|\det(A)|\geq |\alpha|^n$;
        \item if $|\det(A)|=|\alpha|^n$ then $A=\alpha I$.
    \end{itemize}
\end{ex}
\begin{ex}[IMC 2025 P8]
    For $A\in \mathbb{M}_n(\mathbb{R})$ let $A^R$ be the matrix obtained by rotating $A$ counter-clockwise by $90^\circ$; for example
    \[\begin{bmatrix}
    1 & 2 & 3 \\
    4 & 5 & 6 \\
    7 & 8 & 9
    \end{bmatrix}^{R}
    =
    \begin{bmatrix}
    3 & 6 & 9 \\
    2 & 5 & 8 \\
    1 & 4 & 7
    \end{bmatrix}.\]
    Prove that if $A=A^R$ then every eigenvalue $\lambda$ of $A$ satisfies $\mathrm{Re}(\lambda)=0$ or $\mathrm{Im}(\lambda)=0$.
\end{ex}

\newpage
\section{Matrix Norms}
\begin{definition}
    A \emph{norm} on a $\mathbb{C}$-vector space $X$ is a function $\norm{\cdot} : X\to \mathbb{R}$ such that for all $x,y\in X$ and $c\in\mathbb{C}$:
    \begin{itemize}
        \item $\norm{x}\geq 0$, and $\norm{x}=0\iff x=0$;
        \item $\norm{cx}=|c|\norm{x}$;
        \item $\norm{x+y}\leq \norm{x}+\norm{y}$.
    \end{itemize}
    A space with a norm is called \emph{normed}. A norm on a space of matrices is \emph{submultiplicative} if $\norm{AB}\leq \norm{A}\norm{B}$ for all $A,B$.
\end{definition}
\begin{prop}
    If $V$ has an inner product $\langle\cdot,\cdot\rangle$, then $\norm{v}=\sqrt{\langle v,v \rangle}$ defines a norm.
\end{prop}
For the standard inner product on $\mathbb{R}^n$ this is the usual Euclidean length of a vector.
\begin{definition}
    Let $X,Y$ be normed spaces over $\mathbb{C}$. The \emph{operator norm} of a linear map $f:X\to Y$ is
    \[\norm{f}=\sup_{\norm{x}_X=1}\norm{f(x)}_Y=\sup_{x\neq 0} \frac{\norm{f(x)}_Y}{\norm{x}_X}.\]
    Operator norms are always submultiplicative.
\end{definition}
\begin{remark}
    If $X$ is finite dimensional, the supremum is attained, i.e.\ it is a maximum.
\end{remark}
\begin{definition}
    On $\mathbb{M}_n(\mathbb{C})$ we use one of two norms:
    \begin{itemize}
        \item the \emph{Frobenius norm} $\norm{A}_F=\sqrt{\Tr(AA^*)}=\big(\sum_{i,j}|a_{i,j}|^2\big)^{1/2}$;
        \item the \emph{spectral norm}, the operator norm of $A:\mathbb{C}^n\to \mathbb{C}^n$ for the Euclidean norm.
    \end{itemize}
    For the spectral norm $\norm{AA^*}=\norm{A^*A}=\norm{A}^2$, and if $A$ is normal then $\norm{A}$ is the largest absolute value of an eigenvalue of $A$.
\end{definition}
\subsection{Exercises}
\begin{ex}
    Prove that every operator norm is submultiplicative.
\end{ex}
\begin{ex}
    Prove that $\norm{A^*A}=\norm{A}^2$ and $\norm{A}=\norm{A^*}$ for the spectral norm.
\end{ex}
\begin{ex}[SEEMOUS 2022 P3]
    Let $A\in \mathbb{M}_n(\mathbb{C})$ and $\alpha\in \mathbb{C}$, both non-zero, with $A^2+(A^*)^2=\alpha A A^*$. Prove that:
    \begin{itemize}
        \item $\alpha\in \mathbb{R}$ and $|\alpha|\leq 2$;
        \item $A$ is normal.
    \end{itemize}
\end{ex}
\begin{ex}
    Let $x_1,\ldots,x_n$ be real numbers. Show that
    \[\sum_{k=1}^n\Big(\frac{x_1+x_2+\cdots +x_k}{k}\Big)^2\leq (3+2\sqrt{2})(x_1^2+x_2^2+ \cdots + x_n^2).\]
\end{ex}

\newpage
\section{General ideas and problems}
Problems involving both $A$ and $A^*$ are usually solved with the spectral theorem, with the Frobenius norm, or with the propositions of the previous sections. We present several problems solved with these ideas.
\begin{example}\ref{S25P3}
    Let $A\in \mathbb{M}_n(\mathbb{C})$ with $A^*A^2=AA^*$. Prove that $A^2=A$.
    \begin{proof}
        Let $x\in \Ker(A)$. Then $AA^*x=A^*A^2x=0$, so $0=x^*AA^*x=\norm{A^*x}^2$ and $A^*x=0$. Thus $\Ker(A) \subseteq \Ker(A^*)$, and since $\Null(A)=\Null(A^*)$ the two kernels are equal. By Proposition \ref{prop:KA} there is a unitary $U$ with
        \[A=U^*\begin{bmatrix}
            P & O\\
            O & O
        \end{bmatrix}U,\qquad P \text{ invertible},\]
        and $P$ satisfies the same equation $P^*P^2=PP^*$. Taking conjugate transposes gives $(P^*)^2P=PP^*$, hence $P^*P^2=(P^*)^2P$. Cancelling the invertible $P^*$ on the left and $P$ on the right gives $P=P^*$, and then the equation reads $P^3=P^2$, so $P=I$. Therefore $A=U^*(I\oplus O)U$ and $A^2=A$.
    \end{proof}
\end{example}
\begin{example}
    Let $n\geq 2$. Find all integers $k\geq1$ for which the following statement holds:
    \[\textup{``If }A\in \mathbb{M}_n(\mathbb{C})\text{ satisfies } A^kA^*=A, \text{ then } A=A^*.\textup{''} \]
\end{example}
\begin{proof}
    First forget that $A$ is a matrix and work with a number $\lambda$ (later it will be an eigenvalue). The equation becomes $\lambda^k\overline{\lambda}=\lambda$. If $\lambda\neq0$ this gives $\lambda^{k-2}|\lambda|^2=1$; taking absolute values, $|\lambda|=1$, and then $\lambda^{k-2}=1$. We want this to force $\lambda$ to be real, i.e.\ $\lambda=\overline{\lambda}$.
    \begin{itemize}
        \item $k=1$: $\lambda^{-1}=1$, so $\lambda=1$ is real.
        \item $k=2$: $\lambda|\lambda|^2=\lambda$ holds for \emph{every} $\lambda$ with $|\lambda|=1$, so we may take $\lambda$ non-real.
        \item $k=3$: $\lambda=1$, real.
        \item $k=4$: $\lambda^2=1$, so $\lambda=\pm1$, real.
        \item $k\geq5$: $\lambda^{k-2}=1$ has non-real solutions.
    \end{itemize}
    So for $k=2$ and $k\geq5$ the matrix $A=\lambda I_n$, with $\lambda$ a non-real solution, is a counterexample.

    Now let $k\in\{1,3,4\}$. If $x\in \Ker(A^*)$ then $Ax=A^kA^*x=0$, so $\Ker(A^*)\subseteq\Ker(A)$ and, comparing dimensions, $\Ker(A^*)=\Ker(A)$. By Proposition \ref{prop:KA} we may assume that $A$ is invertible. Then $A^*=A^{1-k}$ is a polynomial in $A^{-1}$, so it commutes with $A$; thus $A$ is normal, and after a unitary change of basis it is diagonal. Each diagonal entry satisfies the scalar equation, so it is real, and hence $A=A^*$.

    The statement holds exactly for $k\in\{1,3,4\}$.
\end{proof}

\begin{example}
    Let $A,B\in \mathbb{M}_n(\mathbb{C})$ with $AB^2A=AB$. Prove that:
    \begin{itemize}
        \item $(AB)^2=AB$;
        \item $(AB-BA)^3=O$.
    \end{itemize}
    \begin{proof}
        Let $X=AB$ and $Y=BA$. The hypothesis says $XY=X$, i.e.\ $X(Y-I)=O$. By Sylvester's inequality in kernel form, $n=\Null(X(Y-I))\leq\Null(X)+\Null(Y-I)$. Since $\Rank(AB-I)=\Rank(BA-I)$ we have $\Null(Y-I)=\Null(X-I)$, so $n\leq \Null(X)+\Null(X-I)$. As $\Ker(X) \cap \Ker(X-I)=\{0\}$, this gives $\mathbb{C}^n=\Ker(X)\oplus\Ker(X-I)$. So $X$ acts as $0$ on one summand and as the identity on the other, hence $X^2=X$.

        For the second part, using $X^2=X$ and $XY=X$,
        \[(X-Y)^2=X^2-XY-YX+Y^2=Y^2-YX=-Y(X-Y),\]
        and $Y^3=B(AB)(AB)A=BX^2A=BXA=Y^2$. Therefore
        \begin{align*}
            (X-Y)^3&=(X-Y)^2(X-Y)=-Y(X-Y)^2=Y^2(X-Y),\\
            (X-Y)^4&=(X-Y)^2(X-Y)^2=-Y(X-Y)^3=-Y^3(X-Y)=-Y^2(X-Y),
        \end{align*}
        so $(X-Y)^4=-(X-Y)^3$. Hence the minimal polynomial of $X-Y$ divides $t^4+t^3=t^3(t+1)$, and the eigenvalues of $X-Y$ lie in $\{0,-1\}$. Since $\Tr(X-Y)=\Tr(AB-BA)=0$, the eigenvalue $-1$ cannot occur. So the minimal polynomial divides $t^3$, i.e.\ $(AB-BA)^3=O$.
    \end{proof}
    \begin{remark}
        When we introduce new matrices $X,Y$ in a problem, we may lose sight of relations between them that are not obvious from the new notation: here the nullity of $Y-I$, the identity $Y^3=Y^2$, and the trace of $X-Y$. Be wary when substituting!
    \end{remark}
\end{example}
\begin{example}[Gazeta Matematic\u{a} 498]
    Let $A,B\in \mathbb{M}_n(\mathbb{C})$ with
    \[A^2-B^2-I_n=\frac{1}{3}(AB-BA).\]
    Prove that:
    \begin{itemize}
        \item $\det(A^2-B^2)=\det(A-B)\det(A+B)=1$;
        \item $(AB-BA)^n=O$.
    \end{itemize}
\end{example}
\begin{proof}
    The determinants in the statement suggest looking at
    \[(A-B)(A+B)=A^2-B^2+(AB-BA),\qquad (A+B)(A-B)=A^2-B^2-(AB-BA).\]
    Set $X=A-B$, $Y=A+B$ and $C=AB-BA$. Using the hypothesis $A^2-B^2=I+\frac13C$,
    \[XY=I+\tfrac{4}{3}C,\qquad YX=I-\tfrac{2}{3}C.\]
    Eliminating $C$ gives $XY+2YX=3I$, and also $XY-YX=2C$ and $A^2-B^2=\frac12(XY+YX)$.

    Now use that $XY$ and $YX$ have the same eigenvalues. If $YXv=\lambda v$ with $v\neq0$, then $XYv=(3-2\lambda)v$, so $3-2\lambda$ is an eigenvalue of $XY$, hence of $YX$. Iterating, $1+(-2)^m(\lambda-1)$ is an eigenvalue of $YX$ for every $m$; since there are only finitely many eigenvalues, $\lambda=1$. So all eigenvalues of $YX$ (and of $XY$) are $1$, and $\det(A-B)\det(A+B)=\det(XY)=1$.

    Since $XY=3I-2YX$ is a polynomial in $YX$, a basis that makes $YX$ upper triangular (e.g.\ its Jordan basis) also makes $XY$ upper triangular, both with $1$'s on the diagonal. For instance, if
    \[YX= \begin{bmatrix}
    1 & 1 & 0 & 0 & 0 \\
    0 & 1 & 0 & 0 & 0 \\
    0 & 0 & 1 & 1 & 0 \\
    0 & 0 & 0 & 1 & 1 \\
    0 & 0 & 0 & 0 & 1
    \end{bmatrix}\quad\text{then}\quad XY= \begin{bmatrix}
    1 & -2 & 0 & 0 & 0 \\
    0 & 1 & 0 & 0 & 0 \\
    0 & 0 & 1 & -2 & 0 \\
    0 & 0 & 0 & 1 & -2 \\
    0 & 0 & 0 & 0 & 1
    \end{bmatrix}.\]
    In that basis $A^2-B^2=\frac12(XY+YX)$ is upper triangular with $1$'s on the diagonal, so $\det(A^2-B^2)=1$; and $C=\frac12(XY-YX)$ is strictly upper triangular, hence nilpotent, so $C^n=O$.
\end{proof}
\begin{example}
    Let $A\in \mathbb{M}_n(\mathbb{C})$ with $(AA^*)^2=A^*A$. Show that $A^*A=AA^*$.
\end{example}
\begin{proof}
    Let $x\in \Ker(A)$. Then $(AA^*)^2x=A^*Ax=0$, so $0=x^*(AA^*)(AA^*)x=\norm{AA^*x}^2$ and $AA^*x=0$; by the same trick $0=x^*AA^*x=\norm{A^*x}^2$, so $A^*x=0$. Hence $\Ker(A)\subseteq\Ker(A^*)$, the kernels are equal, and by Proposition \ref{prop:KA} we may assume that $A$ is invertible.

    Now $AA^*$ and $A^*A$ are Hermitian with positive eigenvalues, and they have the same eigenvalues. The equation $(AA^*)^2=A^*A$ shows that if $\mu$ is an eigenvalue of $AA^*$ then $\mu^2$ is an eigenvalue of $A^*A$, hence of $AA^*$. A finite set of positive numbers closed under squaring is $\{1\}$, so all eigenvalues of $AA^*$ equal $1$. Since $AA^*$ is unitarily diagonalizable, $AA^*=I$, and then $A^*A=(AA^*)^2=I=AA^*$.
\end{proof}
\begin{example}\ref{P}
    Since $\det(A^p+I)$ is zero we get that we must have an eigenvalue such that $\zeta^p+1=$ but $\zeta+1\neq 0$. Since the matrix is over $\mathbb{Q}$
     both it's characteristic and minimal polynomials are over $\mathbb{Q}$ considering the dimension is $p$ we argue that $\Phi_{2p}(x)=\Phi_p(-x)| m_A(x)$ so this implies that amongst the $p$ possible eigenvalues of $A$
      $p-1$ of them are the non-real roots of $\Phi_p(-x)$, denote as $\lambda$ the other eigenvalue of $A$, we know that $\tr(A)=\lambda + \sum\limits_{\Phi_p(\zeta)=0} \zeta=\lambda+1$ by Vieta's formulas. All that is left is to invoke the spectral 
      mapping theorem to argue that $\lambda + 1$ is an eigenvalue of $A+I$.
\end{example}
\begin{example}
    For $\alpha$ being real just take the conjugate transpose on both sides,
    taking the spectral norm on $A^2+(A^*)^2=\alpha AA^*$ we get that $|\alpha|\norm{A}^2=\norm{A^2+(A^*)^2}\leq \norm{A^2}+\norm{(A^*)^2}\leq 2\norm{A}^2 $ so $|\alpha|\leq 2$. For the second part take $x\in \ker(A)$ then $Ax=0$ and $x^*A^*=0$ by application of $x$ and $x^*$ on our equation we get 
    $A^*x=0$ so we can pass to an invrtable matrix. Multiplying by the appropriate inverses we get $A(A^*)^{-1}+(A)^{-1}A^*=\alpha I$, set $B= A^{-1}(A^*)$ this turns to $B+B^*=\alpha I$, this means that $B$ is normal and so we can pass to a diagonal.
    The normality of $A$ is equivilant to $BB^*=I$ so we will prove that $|\lambda|^2=1$. Let $v$ such that $Bv=\lambda v$ so $A^*v=\lambda Av$ substituting to the original equation we get $A^2(1+\lambda^2)v=A^2\alpha \lambda v$ since $A^2$ is invertible we get $1+\lambda^2= \alpha \lambda$ meaning that $\lambda=\frac{a\pm \sqrt{a^2-4}}{2}$ in either case $\lambda \overline{\lambda}=1$ so $BB^*=I$.
\end{example}
\newpage

\section{Appendix}
\subsection{From linear functions to matrices}
Let $V,W$ be spaces of finite dimension over some $\mathbb{F}$ and $f:V\to W$ a linear function.
\begin{definition}
    An \emph{ordered basis} of $V$ is a basis together with an ordering of its elements, written as a tuple $\mathcal{B}=(v_1,\ldots,v_n)$ (we care about the order in which we write out the basis). Every $v\in V$ can be written uniquely as $v=\sum_{i=1}^n a_iv_i$, and the \emph{coordinate vector} of $v$ is
    \[[v]_{\mathcal{B}}=\begin{bmatrix} a_1\\ \vdots\\ a_n\end{bmatrix}\in\mathbb{M}_{n\times 1}(\mathbb{F}).\]
\end{definition}
\begin{prop}
    The coordinate map $[\cdot]_{\mathcal{B}}:V\to\mathbb{M}_{n\times1}(\mathbb{F})$ is an isomorphism. It sends $v_i$ to the column with $1$ in the $i$-th place and $0$ elsewhere.
\end{prop}
Let $\mathcal{B}^V=(v_1,\ldots ,v_n)$ and $\mathcal{B}^W=(w_1,\ldots ,w_m)$ be ordered bases of $V$ and $W$.
\[\begin{tikzcd}
	V & W \\
	{\mathbb{M}_{n\times1}} & {\mathbb{M}_{m\times1}}
	\arrow["f", from=1-1, to=1-2]
	\arrow["{[\cdot]_{\mathcal{B}^V}}"', from=1-1, to=2-1]
	\arrow["{[\cdot]_{\mathcal{B}^W}}", from=1-2, to=2-2]
	\arrow["A"', dashed, from=2-1, to=2-2]
\end{tikzcd}\]
The matrix $A$ will be the matrix representation of our function $f$ with respect to $\mathcal{B}^V,\mathcal{B}^W$, i.e.\ the unique matrix with $[f(v)]_{\mathcal{B}^W}=A[v]_{\mathcal{B}^V}$ for all $v$.
Its definition is straightforward because all we have to do is to express each $f(v_i)=\sum_j a_{j,i} w_j$, this will lead us to define the matrix $(A)_{i,j}=a_{i,j}$; in other words the $i$-th column of $A$ is $[f(v_i)]_{\mathcal{B}^W}$.
As an exercise check why the multiplication is compatible with the composition.
With the above idea we can define a map from $\mathcal{L}(V,W)$ to $\mathbb{M}_{m\times n}(\mathbb{F})$ that we will call $M$, this is linear, if we happen to work on $\mathcal{L}(V,V)$ with some fixed basis this map respects compositions, meaning $M(f\circ g)=M(f)M(g)$.
\begin{definition}
    Let $\mathcal{B}=(e_1,\ldots,e_n)$ and $\mathcal{B}'=(e'_1,\ldots,e'_n)$ be two ordered bases of $V$. The \emph{change of basis matrix} from $\mathcal{B}$ to $\mathcal{B}'$ is the matrix $P$ whose $j$-th column is $[e'_j]_{\mathcal{B}}$, i.e.\ $e'_j=\sum_iP_{ij}e_i$. It is invertible and
    \[[v]_{\mathcal{B}}=P\,[v]_{\mathcal{B}'}\qquad\text{for all } v\in V,\]
    so $P$ is the matrix of $\mathrm{Id}_V$ with respect to $\mathcal{B}',\mathcal{B}$.
\end{definition}
Change of basis is the same thing as above: take $f$ to be the identity, since we do not want to alter our input, we just want to express it in a different basis. If we want to change the bases we can consult the following \[\begin{tikzcd}
	V & V & W & W \\
	{\mathbb{M}_{n\times1}} & {\mathbb{M}_{n\times1}} & {\mathbb{M}_{m\times1}} & {\mathbb{M}_{m\times1}}
	\arrow["{Id_V}", from=1-1, to=1-2]
	\arrow["{[\cdot]_{\mathcal{B}^V_2}}"', from=1-1, to=2-1]
	\arrow["f", from=1-2, to=1-3]
	\arrow["{[\cdot]_{\mathcal{B}^V_1}}"', from=1-2, to=2-2]
	\arrow["{Id_W}", from=1-3, to=1-4]
	\arrow["{[\cdot]_{\mathcal{B}^W_1}}", from=1-3, to=2-3]
	\arrow["{[\cdot]_{\mathcal{B}^W_2}}", from=1-4, to=2-4]
	\arrow["X"', dashed, from=2-1, to=2-2]
	\arrow["A"', dashed, from=2-2, to=2-3]
	\arrow["Y"', dashed, from=2-3, to=2-4]
\end{tikzcd}\]
So we get that the matrix that behaves like $f$ with respect to $\mathcal{B}^V_2$, $\mathcal{B}^W_2$ is given by $YAX$.
The fact that $X,Y$ are invertible follows from the fact that the identity is an isomorphism. In terms of change of basis matrices, if $P$ is the change of basis matrix from $\mathcal{B}^V_1$ to $\mathcal{B}^V_2$ and $Q$ the one from $\mathcal{B}^W_1$ to $\mathcal{B}^W_2$, then $X=P$ and $Y=Q^{-1}$, so the new matrix is $Q^{-1}AP$. For $V=W$ and the same basis on both sides this is $P^{-1}AP$, which is exactly similarity of matrices.
\subsection{Inner Products, Gramians and Duality}
Let $V$ be a complex vector space of dimension $n$ with an inner product $\langle\cdot,\cdot\rangle$, linear in the first argument (as in Section 1). Everything below also works over $\mathbb{R}$ by dropping the conjugations.

\begin{definition}
    If $\mathcal{B}=(x_1,x_2,\ldots,x_n)$ is an ordered basis, we define the \emph{Gramian Matrix} $\mathscr{G}=(\langle x_j,x_i \rangle)_{i,j}$, a hermitian matrix.
\end{definition}
If we wish to find $\langle v,w\rangle$ we may write it out as a sum $\sum_{i,j}c_{i,j}\langle x_i,x_j \rangle$, without much trouble you can convince yourself that this is exactly the same as $[w]_{\mathcal{B}}^* \mathscr{G}\,[v]_{\mathcal{B}}$.
Notice that this gives us a quick way of calculating inner products with respect to some change of basis.

\begin{prop}
    Let $\mathscr{G}$ be the Gramian of $\mathcal{B}$. Then
    \begin{itemize}
        \item $\langle v,w\rangle=[w]_{\mathcal{B}}^*\,\mathscr{G}\,[v]_{\mathcal{B}}$ for all $v,w\in V$.
        \item $\mathscr{G}$ is hermitian and positive definite: $X^*\mathscr{G}X>0$ for every nonzero column $X$.
        \item If $P$ is the change of basis matrix from $\mathcal{B}_1$ to $\mathcal{B}_2$ and $\mathscr{G}_1,\mathscr{G}_2$ are the respective Gramians, then, since $[v]_{\mathcal{B}_1}=P[v]_{\mathcal{B}_2}$,
        \[\mathscr{G}_2=P^*\mathscr{G}_1P.\]
        \item $\mathcal{B}$ is orthonormal if and only if $\mathscr{G}=I$.
    \end{itemize}
\end{prop}
\subsection{Rings and Ideals}
All rings are commutative with $1$.
 
\begin{definition}
An \emph{integral domain} is a ring with $1\neq0$ and no zero divisors ($ab=0\Rightarrow a=0$ or $b=0$).
A \emph{field} is a ring in which every nonzero element is a unit.
An \emph{ideal} $I\subseteq R$ is an additive subgroup with $RI\subseteq I$; it is \emph{principal} if $I=(a)=aR$.
A proper ideal $M$ is \emph{maximal} if no ideal lies strictly between $M$ and $R$.
\end{definition}
 
\begin{theorem}
For an ideal $I\neq R$: \ $I$ is maximal $\iff$ $R/I$ is a field.
\end{theorem}
 
\subsection*{PIDs and Euclidean domains}
 
\begin{definition}
A \emph{PID} is an integral domain in which every ideal is principal.
A \emph{Euclidean domain} is an integral domain with a map $N:R\setminus\{0\}\to\mathbb{Z}_{\ge0}$ such that
for all $a$ and $b\neq0$ there are $q,r$ with $a=bq+r$ and $r=0$ or $N(r)<N(b)$.
\end{definition}
 
\begin{theorem}
Every Euclidean domain is a PID.
\end{theorem}
\begin{proof}
Given an ideal $I\neq0$, pick $b\in I\setminus\{0\}$ with $N(b)$ minimal. For $a\in I$, write $a=bq+r$;
then $r\in I$, so minimality forces $r=0$. Hence $I=(b)$.
\end{proof}
 
\begin{theorem}[{$\mathbb{F}[x]$ is a PID}]
Let $\mathbb{F}$ be a field. For $f,g\in \mathbb{F}[x]$, $g\neq0$, there are unique $q,r$ with $f=gq+r$, $r=0$ or $\deg r<\deg g$.
Hence $\mathbb{F}[x]$ is Euclidean with $N=\deg$, and so a PID: every nonzero ideal is generated by any of its
elements of least degree.
\end{theorem}
\begin{proof}[Sketch]
If $\deg f\ge\deg g$, subtract $a_nb_m^{-1}x^{n-m}g$ from $f$ to lower the degree, and induct.
This uses that the leading coefficient $b_m$ of $g$ is invertible.
\end{proof}
 
\begin{remark}
For $A\in \mathbb{M}_n(\mathbb{F})$, the set $\{p\in \mathbb{F}[x]: p(A)=0\}$ is a nonzero ideal (Cayley--Hamilton), so it equals
$(m_A)$ for a unique monic $m_A$, the \emph{minimal polynomial}. The PID property fails for $\mathbb{Z}[x]$:
the ideal $(2,x)$ is not principal.
\end{remark}
 
\begin{corollary}
For nonzero $f\in \mathbb{F}[x]$: \ $f$ irreducible $\iff$ $(f)$ maximal $\iff$ $\mathbb{F}[x]/(f)$ is a field.
\end{corollary}
 
\subsection*{GCD of polynomials}
 
\begin{definition}
For $f,g\in \mathbb{F}[x]$ not both zero, $\gcd(f,g)$ is the monic generator of $(f,g)$. In particular
(B\'ezout) $uf+vg=\gcd(f,g)$ for some $u,v\in \mathbb{F}[x]$. It is computed by the Euclidean algorithm,
so it does not change when $\mathbb{F}$ is replaced by a larger field.
\end{definition}
 
\begin{theorem}
Let $\mathbb{F}\subseteq\mathbb{C}$ and $f,g\in \mathbb{F}[x]$, not both zero. Then
\[
\gcd(f,g)=1 \iff f \text{ and } g \text{ have no common root in } \mathbb{C}.
\]
\end{theorem}
\begin{proof}
If $uf+vg=1$, a common root $\alpha$ would give $0=1$. Conversely, if $\deg\gcd(f,g)\ge1$, it has a
root $\alpha\in\mathbb{C}$, which is then a root of both $f$ and $g$.
\end{proof}
 
\begin{corollary}
$f$ has no repeated root in $\mathbb{C}$ $\iff$ $\gcd(f,f')=1$. In particular, irreducible polynomials over
$\mathbb{Q}$ have distinct roots. Also $\gcd(x^m-1,x^n-1)=x^{\gcd(m,n)}-1$.
\end{corollary}
 
\subsection*{Irreducibility over $\mathbb{Q}$}
 
\begin{theorem}
Let $f=a_nx^n+\dots+a_0\in\mathbb{Z}[x]$.
\begin{enumerate}
  \item \emph{Gauss:} if $f$ factors in $\mathbb{Q}[x]$, it factors in $\mathbb{Z}[x]$ with factors of the same degrees.
  \item \emph{Rational roots:} a root $p/q$ in lowest terms has $p\mid a_0$, $q\mid a_n$.
        For $\deg f\le3$, irreducible $\iff$ no rational root.
  \item \emph{Eisenstein:} if a prime $p$ divides $a_0,\dots,a_{n-1}$, $p\nmid a_n$, $p^2\nmid a_0$, then $f$ is irreducible.
  \item \emph{Mod $p$:} if $p\nmid a_n$ and $\bar f$ is irreducible in $\mathbb{F}_p[x]$, then $f$ is irreducible.
  \item \emph{Shift:} $f(x)$ is irreducible $\iff$ $f(x+c)$ is irreducible.
\end{enumerate}
\end{theorem}
 
\begin{example}
For $p$ prime, $\Phi_p(x+1)=\frac{(x+1)^p-1}{x}=\sum_{k=1}^{p}\binom{p}{k}x^{k-1}$ is Eisenstein at $p$,
so $\Phi_p=x^{p-1}+\dots+x+1$ is irreducible.
\end{example}
 
\subsection*{Cyclotomic polynomials}
 
\begin{definition}
$\displaystyle\Phi_n(x)=\prod_{\substack{1\le k\le n\\ \gcd(k,n)=1}}\bigl(x-e^{2\pi ik/n}\bigr)$, \quad $\deg\Phi_n=\varphi(n)$.
\end{definition}
 
The roots of $\Phi_n$ are the \emph{primitive} $n$-th roots of unity, i.e.\ those of order exactly $n$.
 
\begin{theorem}
\begin{enumerate}
  \item $x^n-1=\prod_{d\mid n}\Phi_d(x)$, hence $n=\sum_{d\mid n}\varphi(d)$.
  \item $\Phi_n\in\mathbb{Z}[x]$ and is monic (induction on $n$, dividing $x^n-1$ by the monic $\prod_{d\mid n,\,d<n}\Phi_d$).
  \item $\Phi_n(x)=\prod_{d\mid n}\bigl(x^d-1\bigr)^{\mu(n/d)}$ (M\"obius inversion).
  \item $\Phi_n$ is irreducible over $\mathbb{Q}$, so it is the minimal polynomial of every primitive $n$-th root
        of unity.
\end{enumerate}
\end{theorem}
\begin{proof}[Idea of (4)]
Let $\zeta$ be a primitive root with minimal polynomial $f\mid\Phi_n$, and $p\nmid n$ a prime.
If $\zeta^p$ were not a root of $f$, then $f(x)$ and $g(x^p)$ would share the root $\zeta$, where $g$ is
the minimal polynomial of $\zeta^p$; so $f\mid g(x^p)$. Mod $p$, $g(x^p)\equiv g(x)^p$, so $\bar f$ and $\bar g$ share
a factor, and $x^n-1$ would have a repeated root in $\overline{\mathbb{F}}_p$. This is impossible, since
$\gcd(x^n-1,nx^{n-1})=1$ in $\mathbb{F}_p[x]$. So $f(\zeta^p)=0$ for all $p\nmid n$; iterating gives $f(\zeta^k)=0$ for all
$\gcd(k,n)=1$, so $f=\Phi_n$.
\end{proof}
 
\begin{prop}[Computing rules]
Let $p$ be prime and $n\ge1$.
\begin{itemize}
  \item $\Phi_p(x)=x^{p-1}+\dots+x+1$ \ and \ $\Phi_{p^k}(x)=\Phi_p\bigl(x^{p^{k-1}}\bigr)$.
  \item $\Phi_{pn}(x)=\Phi_n(x^p)$ if $p\mid n$, \ and \ $\Phi_{pn}(x)=\Phi_n(x^p)/\Phi_n(x)$ if $p\nmid n$.
  \item $\Phi_{2n}(x)=\Phi_n(-x)$ for odd $n>1$.
  \item For $n\ge2$: $\Phi_n(0)=1$ and $\Phi_n$ is palindromic, $x^{\varphi(n)}\Phi_n(1/x)=\Phi_n(x)$.
  \item $\Phi_n(1)=p$ if $n=p^k$, and $\Phi_n(1)=1$ if $n$ has two or more distinct prime factors.
\end{itemize}
\end{prop}
 
\begin{example}
$\Phi_{12}(x)=\dfrac{(x^{12}-1)(x^2-1)}{(x^6-1)(x^4-1)}=\dfrac{x^6+1}{x^2+1}=x^4-x^2+1$,
or, by the rules, $\Phi_{12}(x)=\Phi_6(x^2)=x^4-x^2+1$.
\end{example}
 
\paragraph{Link to matrices.}
If $A\in\mathbb{M}_k(\mathbb{F})$ satisfies $A^n=I$, then $m_A\mid x^n-1$.
\begin{itemize}
  \item Over $\mathbb{C}$: $x^n-1$ has distinct roots, so $A$ is diagonalizable with eigenvalues $n$-th roots of unity.
  \item Over $\mathbb{Q}$: $m_A$ is a product of distinct $\Phi_d$ with $d\mid n$, and the characteristic
        polynomial is a product of (possibly repeated) $\Phi_d$, $d\mid n$.
  \item The order of $A$ is $\operatorname{lcm}\{d:\Phi_d\mid m_A\}$. For example, if $A^p=I$ with $p$ prime
        and $A\neq I$, then $\Phi_p\mid m_A$, so $p-1\le k$.
\end{itemize}
 
\begin{center}
\renewcommand{\arraystretch}{1.15}
\begin{tabular}{>{$}r<{$} >{$}l<{$} @{\qquad} >{$}r<{$} >{$}l<{$}}
\toprule
n & \Phi_n(x) & n & \Phi_n(x)\\
\midrule
1  & x-1               & 11 & x^{10}+x^9+\dots+x+1\\
2  & x+1               & 12 & x^4-x^2+1\\
3  & x^2+x+1           & 14 & x^6-x^5+x^4-x^3+x^2-x+1\\
4  & x^2+1             & 15 & x^8-x^7+x^5-x^4+x^3-x+1\\
5  & x^4+x^3+x^2+x+1   & 16 & x^8+1\\
6  & x^2-x+1           & 18 & x^6-x^3+1\\
7  & x^6+x^5+\dots+x+1 & 20 & x^8-x^6+x^4-x^2+1\\
8  & x^4+1             & 22 & x^{10}-x^9+x^8-\dots-x+1\\
9  & x^6+x^3+1         & 24 & x^8-x^4+1\\
10 & x^4-x^3+x^2-x+1   & 30 & x^8+x^7-x^5-x^4-x^3+x+1\\
\bottomrule
\end{tabular}
 
\smallskip
{\small All $n$ with $\varphi(n)\le10$.}
\end{center}
 

\addcontentsline{toc}{section}{References}
\begin{bibdiv}
\begin{biblist}
\bib{GA17}{book}{
	label={PaB},
      author={Gelca, Razvan},
      author={Andreescu, Titu},
       title={Putnam and Beyond},
     edition={2},
   publisher={Springer International Publishing},
     address={Cham, Switzerland},
      date={2017},
      isbn={978-3-319-58986-2},
       doi={10.1007/978-3-319-58988-6},
}
\bib{Z3}{book}{
   author={Zhang, Fuzhen},
   title={Problems in Linear Algebra and Matrix Theory},
   publisher={World Scientific},
   address={Singapore},
   date={2021},
   isbn={978-981-123-910-6},
   doi={10.1142/12301}
}
\bib{Z}{book}{
   author={Zhang, Fuzhen},
   title={Linear Algebra: Challenging Problems for Students},
   series={The Johns Hopkins Studies in the Mathematical Sciences},
   publisher={The Johns Hopkins University Press},
   address={Baltimore, MD},
   date={1996},
   isbn={978-0-8018-5324-1}
}
\bib{Z2}{book}{
   author={Zhang, Fuzhen},
   title={Matrix Theory: Basic Results and Techniques},
   edition={2},
   series={Universitext},
   publisher={Springer},
   address={New York},
   date={2011},
   isbn={978-1-4614-1098-0},
   doi={10.1007/978-1-4614-1099-7}
}
\bib{Pop2020}{article}{
      author={Pop, Vasile},
       title={Relations between ranks of matrix polynomials},
        date={2020},
     journal={arXiv preprint arXiv:2010.00634},
      eprint={2010.00634},
         url={https://arxiv.org/abs/2010.00634},
}

\bib{Berger1}{book}{
    author={Berger, Marcel},
    title={Geometry I},
    series={Universitext},
    publisher={Springer},
    address={Berlin},
    date={1987},
}
\bib{Berger2}{book}{
    author={Berger, Marcel},
    title={Geometry II},
    series={Universitext},
    publisher={Springer},
    address={Berlin},
    date={1987},
}
\bib{Kaplansky}{book}{
    author={Kaplansky, Irving},
    title={Linear Algebra and Geometry: A Second Course},
    publisher={Dover},
    address={New York},
    date={2003},
}
\bib{lang2002algebra}{book}{
   author={Lang, Serge},
   title={Algebra},
   edition={3},
   series={Graduate Texts in Mathematics},
   volume={211},
   publisher={Springer-Verlag},
   address={New York},
   date={2002},
   isbn={978-0-387-95385-4},
   doi={10.1007/978-1-4613-0041-0}
}
\bib{varsos2012algebra}{book}{
   author={Varsos, Dimitrios},
   author={Deriziotis, Dimitrios},
   author={Emmanouil, Ioannis},
   author={Maliakas, Michalis},
   author={Talelli, Olympia},
   title={Mia Eisagogi stin Algebra},
   language={Greek},
   edition={3},
   publisher={Sofia Publications},
   address={Thessaloniki},
   date={2012},
   isbn={978-960-6706-37-0}
}
\bib{T}{article}{
   author={Thome, N{\'e}stor},
   title={Inequalities and equalities for $l = 2$ (Sylvester), $l = 3$ (Frobenius), and $l > 3$ matrices},
   journal={Aequationes Mathematicae},
   volume={90},
   number={5},
   date={2016},
   pages={951--960},
   doi={10.1007/s00010-016-0412-4},
}
\end{biblist}
\end{bibdiv}
\end{document}
%% Mind !! Add someting about irreducability.